
%%FILE blueprint/src/chapter/Bachlorarbeit/Abstract.tex
\section*{Abstract}
The Bachelor's thesis addresses the formalisation of ancient construction problems, with a particular focus on the "Impossibility of Trisecting the Angle and Doubling the Cube". 
This is the eighth theorem published by Freek Wiedijk in the '100 Theorems' compendium, which serves to allow for a comparison of different theorem provers. 
For this, the theorem prover Lean is utilised. As a result, this paper is devoted to the problems and solutions that arise when formalising ancient construction problems and proves the impossibility of doing so using Galois' theory.
\section*{Zusammenfassung in deutscher Sprache}
Die Bachelorarbeit befasst sich mit der Formalisierung antiker Konstruktionsprobleme, wobei der Schwerpunkt auf der „Unmöglichkeit der Dreiteilung des Winkels und der Verdoppelung des Würfels“ liegt. 
Dies ist das achte Theorem, das Freek Wiedijk in dem im Rahmen von „100 Theorems“ veröffentlichten Kompendium veröffentlicht hat, das dazu dient, einen Vergleich verschiedener Theorembeweiser zu ermöglichen. 
Hierfür wird der Theorembeweiser Lean verwendet. Die vorliegende Arbeit widmet sich daher den Problemen und Lösungen, die bei der Formalisierung alter Konstruktionsprobleme auftreten, und beweist die Unmöglichkeit, dies mit der Galois-Theorie zu tun.

%%FILE blueprint/src/chapter/Bachlorarbeit/Acknowledgment.tex
\section*{Acknowledgment}
I would like to express my gratitude to my advisor, Prof. Dr. Floris van Doorn, for his help during all those times when Lean and I had a different understanding of what a sufficient proof looks like.
Additionally, I am grateful for the encouragement to maintain a broader perspective and to focus on tangible outcomes, as otherwise I would still be engaged in the process of defining a line.

I would like to express my gratitude to Prof. Jan Schröer for providing the impetus for me to attempt to formalise this topic through his stimulating lecture on the subject.

Furthermore, I would like to express my gratitude to my friend Joshua Hendriks, who had to endure an initial, imperfect iteration of this thesis and identify areas for improvement. 
%Furthermore, I would like to express my gratitude to my friends Joshua Hendriks and ???, who had to endure an initial, imperfect iteration of this thesis and identify areas for improvement. 

Finally, I would like to thank Leo Diedering for giving me an idea on how to structure my thesis.

%%FILE blueprint/src/chapter/Bachlorarbeit/BA_Titelseite_LaTeX/meineArbeit.tex
\documentclass[11pt, a4paper, german]{article}

\usepackage[german]{babel}
\usepackage{BA_Titelseite}


%Namen des Verfassers der Arbeit
\author{X Y}
%Geburtsdatum des Verfassers
\geburtsdatum{1. April 1900}
%Gebortsort des Verfassers
\geburtsort{Bonn}
%Datum der Abgabe der Arbeit
\date{\today}

%Name des Betreuers
% z.B.: Prof. Dr. Peter Koepke
\betreuer{Betreuer: Prof. Dr. X Y}
%Name des Zweitgutachters
\zweitgutachter{Zweitgutachter: Prof. Dr. X Y}
%Name des Instituts an dem der Betreuer der Arbeit tätig ist.
%z.B.: Mathematisches Institut
\institut{Institut XYZ}
%\institut{Mathematisches Institut}
%\institut{Institut f\"ur Angewandte Mathematik}
%\institut{Institut f\"ur Numerische Simulation}
%\institut{Forschungsinstitut f\"ur Diskrete Mathematik}
%Titel der Bachelorarbeit
\title{This is only an example}
%Do not change!
\ausarbeitungstyp{Bachelorarbeit Mathematik}

\usepackage{amsmath}
\usepackage{amssymb}
\usepackage{amsthm}
\theoremstyle{plain}
\newtheorem{theorem}{Satz}
\newtheorem{proposition}[theorem]{Proposition}
\newtheorem{lemma}[theorem]{Lemma}
\newtheorem{corollary}[theorem]{Korollar}
\theoremstyle{remark}
\newtheorem{remark}[theorem]{Bemerkung}
\theoremstyle{definition}
\newtheorem{definition}[theorem]{Definition}
\numberwithin{equation}{section}
\numberwithin{theorem}{section}

\usepackage{mathtools}
\DeclarePairedDelimiter\abs{\lvert}{\rvert}
\DeclarePairedDelimiter\norm{\lVert}{\rVert}

\usepackage[T1]{fontenc}
\usepackage[utf8]{inputenc}
\usepackage[hidelinks]{hyperref}
\usepackage[style=alphabetic]{biblatex}
\addbibresource{Literatur.bib}

\begin{document}

\maketitle

\section*{Zusammenfassung}
Wir führen die Benutzung einiger \LaTeX-Befehle und -Pakete anhand von Beispielen vor.
Die Verwendung der hier empfohlenen Lösungen ist, mit Ausnahme des Titelblatts, nicht verpflichtend, sollte Ihre Arbeit aber erheblich vereinfachen.

Dieses Dokument spiegelt den Stand von 2021 wider; wenn Sie es im Jahr 2030 oder später noch lesen, sollten Sie sich vermutlich nach Alternativen für hier vorgestellte Lösungen umschauen.

\begin{flushright}
---Pavel Zorin-Kranich
\end{flushright}

\clearpage
\tableofcontents

\clearpage
\section{Einrichtung Ihres \TeX-Systems}
\subsection{Installation}
Sie benötigen eine aktuelle Version von \TeX Live.
Alle gängigen Linux-Distri\-bu\-ti\-onen stellen dafür Pakete bereit.
Für andere Betriebssysteme ist \TeX Live unter
\url{https://www.tug.org/texlive/}
erhältlich.
Bei der Installation von \TeX Live lassen sich verschiedene Bestandteile auswählen; wenn Sie unsicher sind, was Sie brauchen, können Sie alles installieren.

\subsection{Dokumentation}
Die Distribution \TeX Live beinhaltet Dokumentation für alle darin enthaltenen Pakete.
Eine kurze Einführung in \LaTeX\ sehen Sie, wenn Sie in der Kommandozeile
\begin{verbatim}
$ texdoc lshort-german
\end{verbatim}
eingeben.
Die Dokumentation für ein Paket können Sie sich mit
\begin{verbatim}
$ texdoc PAKETNAME
\end{verbatim}
anzeigen lassen, also z.B.\ \verb|texdoc tikz|.

Wenn Sie ausgefallene mathematische Symbole suchen (die in der o.g.\ Kurzbeschreibung nicht auftauchen), werden Sie hier fündig:
\begin{verbatim}
$ texdoc comprehensive
\end{verbatim}
Es gibt viele Online-Ressourcen zu \LaTeX.
Es kann aber schwierig sein, aktuelle Informationen von veralteten zu unterscheiden.
Die aktivste, und deshalb aktuellste, Ressource ist zur Zeit \url{https://tex.stackexchange.com/}

Es gibt auch Bücher über \LaTeX, wie z.B.\ \cite{More_Math_into_Latex}.

\subsection{Editor}
Wenn Sie noch keinen Lieblingseditor haben, empfehle ich aktuell \emph{Visual Studio Code} mit den Erweiterungen \emph{LaTeX Workshop} (IDE-Funktionen) und \emph{LTeX} (Rechtschreibprüfung).
Spezialisierte \LaTeX-Editoren bieten dem gegenüber keine echten Vorteile, und diesen Editor können Sie auch für viele andere Zwecke einsetzen.

Richten Sie ihren Editor so ein, dass er \verb|pdflatex| und \verb|biber| (und nicht \verb|latex|+\verb|dvipdf| bzw.\ \verb|bibtex|) benutzt.
Qualitäts-Editoren sollten automatisch erkennen, was Sie benötigen.
Wenn Sie keinen Editor mit integriertem Viewer verwenden, sollten Sie noch \verb|pdfsync| manuell einrichten.

\clearpage
\section{Paketempfehlungen}
\label{sec:pakete}
Ich stelle hier kurz die (nach meiner Meinung) für Mathematik wichtigsten Pakete vor.

\subsection{\texttt{amsmath}}
Stellt nummerierte Umgebungen für Gleichungen bereit, wie z.B.
\begin{equation}
\label{eq:1}
F = ma.
\end{equation}
Wenn Sie auf das Display \eqref{eq:1} später verweisen wollen, müssen Sie ihm mit dem Befehl \verb|\label| einen Namen zuweisen, auf den Sie dann mit dem Befehl \verb|\eqref| verweisen können.
Ähnliches (mit \verb|\ref| statt \verb|\eqref|) ist für Kapitel, wie z.B.\ Kapitel~\ref{sec:pakete}, und Sätze, wie z.B.\ Satz~\ref{thm:Cauchy-Schwarz}, möglich.

\subsection{\texttt{amssymb}}
Definiert wichtige Symbole, wie z.B. $\mathbb{R}$ mit dem Befehl \verb|\mathbb{R}|.

\subsection{\texttt{amsthm}}
Stellt den Befehl \verb|\newtheorem| bereit, der benutzt wird um Umgebungen für Sätze, Definitionen, und Ähnliches zu definieren.
Hier einige Beispiele.

\begin{definition}
Ein \emph{metrischer Raum} besteht aus einer Menge $M$ und einer Abbildung $d : M \times M \to \mathbb{R}_{\geq 0}$, genannt \emph{Metrik}, sodass für alle $x,y,z\in M$ gilt
\begin{enumerate}
\item $x=y \iff d(x,y)=0$,
\item $d(x,y)=d(y,x)$ (Symmetrie),
\item $d(x,z) \leq d(x,y)+d(y,z)$ (Subadditivität oder Dreiecksungleichung).
\end{enumerate}
\end{definition}

\begin{theorem}[Cauchy--Schwarz-Ungleichung]
\label{thm:Cauchy-Schwarz}
Sei $H$ ein Hilbertraum und $v,w \in H$.
Dann gilt
\[
\abs{ \langle v,w \rangle }
\leq
\norm{v} \norm{w}.
\]
\end{theorem}
\begin{proof}
Wenn der Beweis mit einer Gleichung endet, können Sie \verb|\qedhere| verwenden, um das Symbol für Beweisende auf der gleichen Zeile anzuzeigen.
\end{proof}
\begin{remark}
Diese Bemerkung ist nicht nötig um den Rest des Textes zu verstehen.
\end{remark}

\begin{proof}[Beweis von Satz~\ref{thm:Cauchy-Schwarz}]
Wenn zwischen der Aussage eines Satzes und dem Beweis noch etwas steht, sollte man explizit sagen, was man gerade beweisen will.
\end{proof}

\subsection{\texttt{fontenc}}
\TeX\ ist ein sehr altes System und benutzt standardmäßig Rasterschriftarten.
Dieses Paket mit der Option \verb|T1| ersetzt sie durch Vektorschriftarten.

\subsection{\texttt{inputenc}}
\TeX\ ist ein sehr altes System und akzeptiert standardmäßig nur ASCII-Eingabe.
Dieses Paket mit der Option \verb|utf8| erlaubt Unicode-Eingabe, bietet also insbesondere Unterstützung für Umlaute und andere Akzente.

\subsection{\texttt{hyperref}}
Erzeugt klickbare Hyperlinks innerhalb des Dokuments bei Benutzung von \verb|\ref| und  \verb|\eqref| und ins Internet mit \verb|\url| und \verb|\href|.

\subsection{\texttt{biblatex}}
Formatiert die Bibliographie aus Daten, die in einer \verb|.bib|-Datei vorliegen.
Die Datei muss mit \verb|\addbibresource| eingebunden werden.
Die Bibliographie wird mit dem Befehl \verb|\printbibliography| gesetzt.
Es ist ein Zwischenschritt mit dem Programm \verb|biber| notwendig, der von Qualitäts-Editoren automatisch ausgeführt wird; Details entnehmen Sie der Paket-Dokumentation.

Bibliographische Daten in dem passenden Format finden Sie in \emph{Math\-Sci\-Net} (\url{https://mathscinet.ams.org/}) oder \emph{zbMATH} (\url{https://zbmath.org/}).
Wenn Sie eine Veröffentlichung, die Sie zitieren wollen, dort nicht finden, wie z.B.\ \cite{arxiv:1507.02515}, können Sie den Eintrag von Hand ausfüllen, oder Sie benutzen einen \verb|.bib|-Manager, wie z.B.\ \emph{KBibTeX} (dieser hat den Vorteil einer eingebauten MathSciNet-Suche).

Wenn Sie aus einem Buch zitieren, benennen Sie immer den relevanten Abschnitt, also z.B.\ \cite[Satz III.2.2]{MR1232192} oder zumindest \cite[p.\ 107]{MR1232192} (die Seitenzahl-Variante ist schlechter, da eher von der Ausgabe abhängig), und nicht einfach \cite{MR1232192}.
In diesem Beispiel ersparen Sie \emph{jedem} Leser Ihrer Arbeit den Aufwand, in einem 700-seitigen Buch etwas Bestimmtes zu suchen.

\clearpage
\section{Optionale Pakete}
Die hier vorgestellten Pakete decken Aufgaben ab, die sich nicht in jeder Arbeit stellen, aber häufig genug vorkommen dürften, um sie an dieser Stelle erwähnenswert zu machen.

\subsection{\texttt{mathtools}}
Vereinfacht häufige Formatierungsaufgaben, wie z.B.\ Einrückung von Gleichungen in der \verb|align|-Umgebung mit dem Befehl \verb|\MoveEqLeft| oder hübsche Formatierung von Klammern, Beträgen, Normen, inneren Produkten, etc.\ mit dem Befehl \verb|\DeclarePairedDelimiter|.

\subsection{\texttt{tikz}}
Das umfangreiche Paket zur Erstellung von zweidimensionalen Grafiken aller Art.

\subsection{\texttt{tikz-cd}}
Eine auf \verb|tikz| aufbauende Lösung zur Erstellung von kommutativen Diagrammen.
Komfortabler und flexibler als alle früheren Lösungen.

Mit \emph{Quiver} (\url{https://q.uiver.app/}) existiert eine WYSIWYG-Ober\-fläche zur Erstellung von \verb|tikz-cd|-Diagrammen.

\subsection{\texttt{graphicx}}
Einbindung externer Grafiken.
Empfehlenswerte Formate sind PDF (für Vektorgrafiken), PNG (für Rastergrafiken), oder JPG (für Fotos).

\subsection{\texttt{booktabs}}
Schönere Tabellen.

\clearpage
\addcontentsline{toc}{section}{Literatur}
\printbibliography

\end{document}


%%FILE blueprint/src/chapter/Bachlorarbeit/Introduction.tex
\chapter[Introduction]{Introduction in the topic of the Bachelor Thesis}

This paper formalises the proof of "The Impossibility of Trisecting the Angle and Doubling the Cube" in ruler and compass construction. 
The Github repository with the code and bluebrint can be found here : \url{https://github.com/Louis-Le-Grand/Formalisation-of-constructable-numbers}
\section{Construction Problems}
The straightedge and compass constructions were developed by the ancient Greeks. 
It consists of an initial set $\mathcal{M}$ of constructed points, a ruler that has no measurements and can draw indefinite lines through at least two existing points,
and a compass, the centre of which is an already constructed point and the radius of which is the distance between two already constructed points. 
To construct new points, take the intersection of two lines, two circles or a line and a circle. The set of all possible intersections is called $\mathcal{M}_{\infty}$. 
In order to construct a line or a circle, it is necessary that our initial set contain at least two points. Consequently, we can normalise our set and assume that $0, 1 \in \mathcal{M}$.

We may now proceed to define the problem of doubling a cube. 
The volume of a cube is equal to the cube of the length of an edge, which may be expressed as $a^3$, where $a$ is the length of an edge. 
Therefore, a cube with a doubled volume, $2\cdot a^3$, has an edge length of the cube root of two times the length of the original edge. 
If we now take the cube with a length of one, the problem is as follows:
\begin{problem}[Doubling the Cube]
    Given the set $\mathcal{M} = \{0, 1\}$, can we construct the point $\sqrt[3]{2}$, i.e. is $\sqrt[3]{2} \in \mathcal{M}_{\infty}$?
\end{problem}

A similar approach allows the problem of trisecting angles to be simplified. 
If we take two points on the real axis and a third point on the unit circle, this defines an angle. Thus we need the points $0$, $1$ and $e^{i\alpha}$. 
The trisection intersects the unit circle at $e^{i\alpha/3}$. Therefore the problem can be described as follows:

\begin{problem}[Trisecting the Angle]
    Given the set $\mathcal{M} = \{0, 1, e^{\imath\cdot\alpha}\}$, can we construct the point $e^{\imath\cdot\frac{\alpha}{3}}$, i.e. is $e^{\imath\cdot\frac{\alpha}{3}} \in \mathcal{M}_{\infty}$?
\end{problem}

\section{Motivation}
The subject of this formalisation is the eighth theorem of Freek Wiedijk's list of "100 Theromes", entitled "The Impossibility of Trisecting the Angle and Doubling the Cube" \cite{Wiedijk}. 
This is a list of theorems based on an online list from 1999 of the 100 most significant theorems in mathematics \cite{Abad_Abad}, which is used for comparative purposes with respect to different theorem provers. 
The list provides a concise overview of the most important theorem provers, showcases fields in which theorems can be used, and also presents some smaller programs that have formalised theorems that had not yet been formalised in any other environment.
 Notably, the aforementioned theorem had not yet been formalised in Lean \cite{Lean_Community}.

\section{Lean, Theorem Prover and Mathlib}
Lean is a functional programming language that can be utilised as an interactive theorem prover. 
The Lean Project was initiated by Leonardo de Moura at Microsoft Research Redmond in 2013 and represents a long-term research endeavour, published under the Apache 2.0 licence.

In order to verify a theorem, it is necessary to find proof. 
Computers can be of assistance in two distinct ways: 
firstly, through interactive theorem proving, which verifies the correctness of a proof step by step; 
and secondly, through automated theorem proving, which attempts to find a proof for a given statement. 
Lean represents a hybrid of these two approaches. 
It is an interactive theorem prover, but it incorporates automated tools and methods to assist in the construction of a fully specified axiomatic proof.\cite{Avigad_deMoura_Kong_Ullrich}


Mathlib is a library of mathematical content for the Lean programming language. It is a community-developed project, created by a large number of contributors and covering a substantial breadth of mathematical topics. To ascertain which areas are encompassed by MathLib, one may refer to the following link: \url{https://leanprover-community.github.io/mathlib-overview.html}. 


\section*{Outline of the Thesis}
The following three chapters present the blueprint for my formalisation.
Consequently, the document includes a number of elementary lemmas and details that are not typically documented in such a manner. 
However, a few statements and most of the auxiliary lemmas are not included in the blueprint.
The structure of this formalisation is based on the lecture "Einführung in die Algebra" by Jan Schröer (WS 22/23).

%%FILE blueprint/src/chapter/Bachlorarbeit/lean.tex
\chapter{A sample of Lean}
The following chapter provides a concise overview of the author's experience with lean, the challenges encountered, and the author's approach to formalisation. 

However, it will not explain the operational logic of lean nor serve as a guide to learning lean, as this is beyond the scope of this thesis.
 For those interested in learning more about lean, I recommend consulting the following resources: \url{https://leanprover-community.github.io/learn.html}

\section{Blueprint}
\begin{wrapfigure}{r}{0.5\textwidth}
    \centering
    \includegraphics[angle=0, width=0.5\textwidth]{Blueprint}
    \label{fig:Blueprint}
    \caption{Blueprint}
    \scriptsize
    \begin{itemize}
        \setlength\itemsep{-0.5em}
        \item[$\checkmark$] Indicates the statment is Formalised
        \item[$\#$] Links to the proof
        \item[\scalebox{0.1}{\includegraphics{gear}}] Shows to the proof
        \item[\scalebox{0.1}{\includegraphics{mindmap}}] Shows dependency
        \item[\scalebox{0.5}{$\text{L}\exists\forall\text{N}$}] Links to the documentation
    \end{itemize}
\end{wrapfigure}

In order to structure my project, I employ the Lean Blueprint tool, created by Patrick Massot. \cite{Massot_2020}
This tool generates a web version of the LaTeX file, which provides an outline of the project and facilitates the integration of code with the associated documentation. 
Please refer to Figure \ref{fig:Blueprint} for an example of the linking. 
Furthermore, it generates a dependency graph \ref{fig:DependencyGraph}, which illustrates the extent of formalisation.
The presented approach is particularly beneficial for larger-scale projects involving multiple contributors, as evidenced in Terence Tao blog post. \cite{Tao_2023}



\clearpage
\begin{figure}[h]
    \centering
    \scalebox{0.5}{\includegraphics[angle=90]{DependencyGraph}}
    \label{fig:DependencyGraph}
    \caption{Dependency Graph}
\end{figure}
\clearpage
\section{A sample of Lean code} 
The following is a representative sample of the Lean code, which is available in its entirety on the GitHub repository of this project. 
\url{https://github.com/Louis-Le-Grand/Formalisation-of-constructable-numbers}

I defined lines as a structure of two points $z_1$ and $z_2$. 
For this structure I then define points by $\{\lambda z_1 + (1-\lambda)z_2\mid\lambda\in\R\}$ which are all the points the line goes through. 
I have omitted the condition $z_1\ne z_2$ because in some settings a line that consists of only one point makes sense.
\begin{lstlisting}
    structure line where
        (z₁ z₂ : ℂ)

    def line.points (l: line) : Set ℂ:= 
        {(t : ℂ) * l.z₁ + (1-t) * l.z₂ | (t : ℝ)}
\end{lstlisting}

To define the circles, I used $c$ of $\C$ as the centre and $r$ as the radius. 
For the points, I could refer to spheres already defined in Mathilb, which had the advantage that I could use existing lemmas about them. 
\begin{lstlisting}
    structure circle where
        (c : ℂ)
        (r : ℝ)

    def circle.points (c: circle) := Metric.sphere c.c c.r
    noncomputable def circle.points' (c: circle) := 
        (⟨c.c, c.r⟩ : EuclideanGeometry.Sphere ℂ)
\end{lstlisting}

To prove that $M_{\infty}$ is a subfile of $\C$, I had to define a new object of type subfile of $\C$, and use $M_{\infty}$ as the subordinate carrier.
This is because each object in Lean is a type, and you cannot switch between types.
\begin{lstlisting}
    noncomputable def MField (M: Set ℂ)(h₀: 0 ∈ M)(h₁: 1∈ M):
            Subfield ℂ where
        carrier := M_inf M
        zero_mem' := by exact M_M_inf M h₀
        one_mem' := by exact M_M_inf M h₁
        add_mem' := by apply add_M_Inf M h₀
        neg_mem' := by apply z_neg_M_inf M h₀
        mul_mem' := by apply mul_M_inf M h₀ h₁
        inv_mem' := by apply inv_M_inf M h₀ h₁
\end{lstlisting}

\section{Conclusion}
The proceeding was an account of the process by which the formalisation of "The Impossibility of Trisecting the Angle and Doubling the Cube" was reached. 
On the one hand, working with lean demonstrates which aspects have not yet been resolved, during the formalisation process and ensures that all outcomes will be accurate. 
However, it requires a more complex proof structure  than is typically required, rendering processes that are otherwise straightforward challenging to demonstrate.
Consequently, this approach is inherently time-consuming. 
This has resulted in the following aspects not being completed in time.

In Lemma \ref{lem:ConjClosed.Intersection_circle_circle}, while calculating the existence of the line through the intersection points of the two circles, it was found that for $\Im (c_1) = \Im (c_2) = 0$ , the lemma is not true.
Due to time limitations, these issues could not be addressed.

In the proof of Theorem \ref{thm:Classfication_z_in_M_inf} the fact that it is a finte adjunction of square roots was not proven and the equivalence of Lemma \ref{lem:Z_in_Minf_imp_eq} is not formalized.

It should be noted that proof of Lemma \ref{lem:dergree_two_eq_sqr} was never forthcoming, as there was a desire to prove the most general setting in order to contribute it to Mathlib.
However, due to the limitations of time, this was not feasible.

Nonetheless, I'm glad that I could make a small contribution to the formalization of the 100 theorems in Lean.

%%FILE blueprint/src/chapter/biblio.tex
\bibliographystyle{plainurl} % We choose the "plain" reference style
\bibliography{construct}




%%FILE blueprint/src/chapter/ch01introduction.tex
\chapter{Basic Constructions with Ruler and Compass}

\section{Defining lines, circles and the set of constructable points}
First we need to define what construction using a ruler and compass means.
We will use $\mathbb{C}$ as plane of drawing and $\mathcal{M} \subset \mathbb{C}$ as the set of constructed points.

\begin{definition}[Line]
    \label{def:line}
    \lean{Construction.line}
    \leanok
    A line $L := (x,y)$ is defined by two points $x,y\in\mathbb{C}$ with $x\ne y$ and a set $$l:=\{\lambda x+(1-\lambda)y\mid\lambda\in\mathbb{R}\}.$$
    We say two lines are equal if their generated set is equal.
\end{definition}
\begin{remark}
    To express lines in the most general way, they are stated without the requirement that $x\ne y$, which allows for trivial lines. 
    The condition is only included in the lemmas that require it.
\end{remark}


\begin{definition}[Circle]
    \label{def:circle}
    \lean{Construction.circle}
    \leanok
    A circle $C := (c,r)$ is defined by a center $c$ with  $c\in\mathbb{C}$, a radius $r\in\mathbb{R}_{\ge 0}$ and the set: $$c:=\{z\in\mathbb{C} \mid\|z-c\|=r\}.$$ 
    We say two circles are equal if their underlying sets are equal.
\end{definition}

\begin{lemma}
    \label{lem:circle_not_eq_iff}
    \lean{Construction.circle_not_eq_iff, Construction.circle_not_eq_iff_radius}
    \leanok
    \uses{def:circel}
    If we have two circles $C_1 = (c_1,r_1)$, $C_2 = (c_2,r_2)$ and we have $c_1 \ne c_2$ or $r_1  \ne r_2 $ then the circles aren't equal. 
\end{lemma}
\begin{proof}
    \leanok
    \uses{def:circle}
    Uses the proof of Eucliden speheres from MathLib.
\end{proof}


\begin{definition}[Set of lines]
    \label{def:set_of_lines}
    \lean{Construction.L, Construction.l_in_L_M}
    \leanok
    \uses{def:line}
    $\mathcal{L(M)}$ is the set of all real straight lines $l$, with $| l\cap \mathcal{M} |\ge 2$. As a set this is:
    \begin{equation*}
        \mathcal{L(M)} := \{l \mid l = \{x,y\} \text{ with }x,y \in \mathcal{M} \land x \neq y\}.
    \end{equation*}
\end{definition}

\begin{definition}[Set of circles]
    \label{def:set_of_circles}
    \lean{Construction.C, Construction.c_in_C_M}
    \leanok
    \uses{def:circle}
    $\mathcal{C(M)}$ is the set of all circles in $\mathbb{C}$, with center in $\mathcal{M}$ and radius of $\mathcal{C}$ which is the distance of two points in $\mathcal{M}$. As an equation this is:
    $$\mathcal{C(M)}:=\{c\mid c = \langle c, \dist r_1 r_2 \rangle \text{ with } c,r_1,r_2\in M\}.$$
\end{definition}

\begin{definition}[Rules to construct a point]
    \label{def:rules_to_constructed_a_point}
    \lean{Construction.ill,Construction.ilc,Construction.icc,Construction.ICL_M}
    \leanok
    \uses{def:set_of_lines, def:set_of_circles}
    We define operations that can be used to construct new points.
    \begin{enumerate}
        \item $(ILL)$ is the intersection of two different lines in $\mathcal{L(M)}$.
        \item $(ILC)$ is the intersection of a line in $\mathcal{L(M)}$ and a circle in $\mathcal{C(M)}$.
        \item $(ICC)$ is the intersection of two different circles in $\mathcal{C(M)}$.
    \end{enumerate}
    $ICL(\mathcal{M})$ is the union of all points that can be constructed using the operations $(ILL)$, $(ILC)$ and $(ICC)$ and $\mathcal{M}$.
\end{definition}

\begin{definition}[Set of constructable points]
    \label{def:set_of_constructable_points}
    \lean{Construction.M_I,Construction.M_inf}
    \leanok
    \uses{def:rules_to_constructed_a_point}
    We inductively define the chain
    \begin{equation*}
        \mathcal{M}_0 \subseteq \mathcal{M}_1 \subseteq \mathcal{M}_2 \subseteq \dots
    \end{equation*}
    with $\mathcal{M}_0 = \mathcal{M}$ and $\mathcal{M}_{n+1} = ICL(\mathcal{M}_n)$\newline
    and call $\mathcal{M}_{\infty} = \bigcup_{n \in \mathbb{N}} \mathcal{M}_n$ the set of all constructable points.
\end{definition}

\begin{remark}
    The set of lines, circles and their intersection $icc$, $ilc$, $icc$ are monoton, i.e. for $N\subseteq M$ is $\mathcal{L(N)}
    \subseteq\mathcal{L(M)}$,\dots .
\end{remark}

\section{Properties of Lines}
In light of the fact that a considerable proportion of the forthcoming project will entail the manipulation of lines and the absence of an existing formalisation from MathLib, 
it is imperative to illustrate the fundamental properties that will be utilised throughout the remainder of the project.

\begin{lemma} 
    \label{lem:line_pionts_iff}
    \lean{Construction.line_pionts_iff}
    \leanok
    \uses{def:line}
    The set of points of a line $l$ is not defined by one position vector, i.e. for every $\alpha\in l$  $\{\lambda x+(1-\lambda)y\mid\lambda\in\mathbb{R}\} = \{\lambda x+\lambda y+\alpha\mid\lambda\in\mathbb{R}\}$.
\end{lemma}
\begin{proof}
    \leanok
    Two sets are equal if they have the same elements. 
    Since $\alpha$ is in $l$, there exists a $\lambda_{\alpha}$ such that $\alpha = \lambda_{\alpha} x+(1-\lambda_{\alpha})y$.
    
    "$\Rightarrow:$" Let $z=\lambda_0 x + (1-\lambda_0)y$ (i.e.  $z\in\{\lambda x+(1-\lambda)y\mid\lambda\in\mathbb{R}\}$), we have to show that $z\in \{\lambda x+\lambda y+\alpha\mid\lambda\in\mathbb{R}\}$,
    which is equivalent to the existence of $\lambda$ such that $z = \lambda x+\lambda y+\alpha$. If we use $\lambda = \lambda_0 - \lambda_{\alpha} $ we get
    $$(\lambda_0 - \lambda_{\alpha}) x + ((\lambda_0 - \lambda_{\alpha}))y + \alpha  \overset{\alpha = \lambda_{\alpha} x+(1-\lambda_{\alpha})y}{=} \lambda_0 x - \lambda_0 y + y = z$$ % \lambda_{\alpha} x+(1-\lambda_{\alpha})y + (\lambda_0 - \lambda_{\alpha}) x + ((\lambda_0 - \lambda_{\alpha}))y
    $\Leftarrow:$ Let $z$ be in $\{\lambda x+(\lambda)y+\alpha\mid\lambda\in\mathbb{R}\}$, if we use $\lambda = \lambda_0 + \lambda_{\alpha} $ we get
    $$(\lambda_0 + \lambda_{\alpha}) x + (1 -(\lambda_0 + \lambda_{\alpha}))y = \lambda_0 x - \lambda_0 y + \lambda_{\alpha} x - \lambda_{\alpha} y + y \overset{\alpha = \lambda_{\alpha} x+(1-\lambda_{\alpha})y}{=}  \lambda_0 (x - y) + \alpha = z$$
\end{proof}

\begin{definition}[parallel]
    \label{def:parallel}
    \lean{Construction.parallel}
    \leanok
    \uses{def:line}
    Two lines $l_1$ and $l_2$ are parallel if they have the same slope, i.e. $\exists r\in\C: l_1 = \{x + r \mid x\in l_2\}$.
\end{definition}

\begin{lemma}[Parallel defined by basis]
    \label{lem:parallel_by_basis}
    \lean{Construction.parallel_basis}
    \leanok
    \uses{def:line, def:parallel}
   If the basis of two lines is moved by the same amount, i.e $l_1.z_1 - l_2.z_1 = l_1.z_2- l_2.z_2$, then they are parallel.
\end{lemma}
\begin{proof}
    \leanok
    \uses{def:line, def:parallel}
    We have to lines $l_1, l_2$ with basis $z_1, z_2$ and $l_1.z_1 - l_2.z_1 = l_1.z_2- l_2.z_2$. To show that they are parallel we have to show that $$\exists r: l_1 = \{z + r \mid z\in l_2\}.$$
    To make it easier to read we define $a := l_1.z_1$, $b:= l_1.z_2$, $x:= l_2.z_1$ and $y:= l_2.z_2$.\\
    Upon unravelling the definition, it becomes evident that we need to show that  $\exists t: t \cdot a +(1-t)\cdot b=z \iff \exists s :s \cdot x+(1-s)y+(a-x)=z$.\\
    $\Rightarrow:$ Claim $s = \frac{t(a-b)}{x-y}$ is a solution.\\
    \begin{proof}
    \begin{align*}
        && z &= s \cdot x+(1-s)y+(a-x)&&| z= t\cdot a+(1-t)\cdot b\\
        &\Leftrightarrow & t\cdot a+(1-t)\cdot b  &= s \cdot x+(1-s)y+(a-x)\\
        &\Leftrightarrow & t\cdot a-t\cdot b +b &= s \cdot x-s\cdot y + y+(a-x)&& | -y\\
        &\Leftrightarrow & t\cdot a-t\cdot b + b-y  &= s \cdot x-s\cdot y+(a-x)-y&&| a-x =b-y\\
        &\Leftrightarrow & t\cdot a-t\cdot b  &= s \cdot x+s\cdot y\\
        &\Leftrightarrow & t(a-b)  &= s(x-y)&&|s:= \frac{t(a-b)}{x-y}\\
        &\Leftrightarrow & t(a-b)  &= \frac{t(a-b)}{x-y}(x-y)\\
        &\Leftrightarrow & t(a-b)  &= t(a-b)
    \end{align*}
    \end{proof}
    "$\Leftarrow:$" Claim $t = \frac{s(x-y)}{a-b}$ is a solution.
    \begin{proof}
    \begin{align*}
        && t\cdot a+(1-t)b &= z&&| z:= s \cdot x+(1-s)y+(a-x)\\
        &\Leftrightarrow& t\cdot a-t\cdot b+b &= s \cdot x-s\cdot y+y+(a-x)&&|-y;a-x =b-y\\
        &\Leftrightarrow& t(a-b) &= s(x-y)&&|t:= \frac{s(x-y)}{a-b}\\
        &\Leftrightarrow& \frac{s(x-y)}{a-b}(a-b) &= s(x-y)
    \end{align*}
    \end{proof}
\end{proof}

\begin{definition}
    \label{def:direction_vector}
    \lean{Construction.direction_vector}
    \leanok
    \uses{def:line}
    The direction vector of a line $l$ is the vector $l.z_1 - l.z_2$.
\end{definition}

\begin{definition}
    \label{def:parallel'}
    \lean{Construction.parallel'}
    \leanok
    \uses{def:line, def:direction_vector}
    Two lines $l_1$ and $l_2$ are \emph{parallel'} if there exists a $k$ such that $l_1.z_1 - l_1.z_2 = k \cdot (l_2.z_1 - l_2.z_2)$.
\end{definition}

\begin{lemma}[Parallel imp parallel']
    \label{lem:parallel_iff_parallel'}
    \lean{Construction.parallel_iff_parallel'}
    \leanok
    \uses{def:parallel, def:parallel'}
    If two lines $l_1$ and $l_2$ are parallel' they are parallel.
\end{lemma}
\begin{remark}
    The other direction holds as well, but is not needed.
\end{remark}

\begin{proof}
    \leanok
    \proves{lem:parallel_iff_parallel'}
    \uses{def:line, def:direction_vector, def:parallel, def:parallel'}
    Let $l_1$ and $l_2$ be two lines with $l_1.z_1 - l_1.z_2 = k \cdot (l_2.z_1 - l_2.z_2)$.
    We have to show that $l_1 = \{z + r \mid z\in l_2\}$, i.e. $\exists r: l_1 = \{z + r \mid z\in l_2\}$.\\
    First we remark that $k\ne 0$ since $l_1.z_1\ne l_1.z_2$. We define $r:= y_1 - y_2$ and show that $l_1 = \{z + r \mid z\in l_2\}$.
    For $z$ to be in $l_1$ is equivalent to $z = \lambda x_1 + (1-\lambda)y_1$ for some $\lambda_0$. 
    For $z$ to be in $\{x + y_1-y_2 \mid z\in l_2\}$  is wquivalent to the existence of $\lambda$ such that 
    $$z = \lambda x_2 + (1-\lambda)y_2 + y_1 - y_2 = \lambda (x_2 - y_2) + y_1 \overset{\lambda = \lambda_0 * k}{=} \lambda_0 k (x_2-x_2)-y_1\overset{x_1-y_1=k(x_2-y_2)}{=}\lambda_0(x_1-y_1)+y_1.$$
\end{proof}

\begin{lemma}[Parallel imp equal]
    \label{lem:eq_of_parallel}
    \lean{Construction.eq_of_parallel}
    \leanok
    \uses{def:line, def:parallel'}
    If two lines $l_1$ and $l_2$ are parallel' and intersect, then they are equal.
\end{lemma}

\begin{proof}
    \leanok
    \uses{def:line, def:parallel', def:direction_vector, lem:line_pionts_iff}
    Let $z\in l_1 \cap l_2$. Then we use lemma \ref{lem:line_pionts_iff} to get $l_1 = l_2 \Leftrightarrow \{\lambda (x_1-y_1)\mid\lambda\in\mathbb{R}\} = \{\lambda (x_2- y_2)+z\mid\lambda\in\mathbb{R}\}$.
    Since for every $\lambda$ we have $\lambda (x_1-y_2) + x = \lambda \cdot k(x_2- y_2)+z$ we get $l_1 = l_2$.
 \end{proof}

\section[The set of constructable points]{The Set $M_{\infty}$}\label{set_of_constructable_points}
In order to work with the set of constructable points, it is necessary to have some basic lemmas at one's disposal. 
While these may appear to be trivial on paper, they are invaluable when undertaking a formalisation.

\begin{lemma}[$\mathcal{M}\subseteq ICL(\mathcal{M})$]
    \label{lem:M_in_ICL_M}
    \lean{Construction.M_in_ICL_M}
    \leanok
    \uses{def:rules_to_constructed_a_point}
    Every set $M$ is included in the constructable points of $M$, i.e.
    \[M \subseteq ICL(M)\]
\end{lemma}
\begin{proof}
    \leanok
    Follows from the definition of $ICL(M)$.
\end{proof}

\begin{lemma}[$\mathcal{M}_i$ Monoton]
    \label{lem:M_i_monoton}
    \lean{Construction.M_I_Monotone,Construction.M_I_Monotone_elm,Construction.M_I_Monotone', Construction.M_I_Monotone_elm'}
    \leanok
    \uses{def:set_of_constructable_points}
    The set $\mathcal{M}_i$ is monoton, i.e. $$\mathcal{M}_i \subseteq \mathcal{M}_{i+1}.$$
    This can be generalised to $$\forall m :\forall n\le m : \mathcal{M}_n \subseteq \mathcal{M}_m.$$
\end{lemma}
\begin{proof}
    \leanok
    Follows from the definition of $\mathcal{M}_i$.
\end{proof}

\begin{lemma}[$\mathcal{M}$ in $\mathcal{M}_i$]
    \label{lem:M_in_Mi}
    \lean{Construction.M_in_M_I }
    \leanok
    \uses{def:set_of_constructable_points}
    The set $\mathcal{M}$ is in $\mathcal{M}_i$, i.e. $$\mathcal{M} \subseteq \mathcal{M}_i.$$
\end{lemma}
\begin{proof}
    \leanok
    \uses{lem:M_i_monoton}
    Combining the fact that $\mathcal{M}_0 = \mathcal{M}$ \ref{def:set_of_constructable_points} and the monotonity of $\mathcal{M}_i$ \ref{lem:M_i_monoton} we get the result.
\end{proof}

\begin{lemma}[$\mathcal{M}_i$ in $\mathcal{M}_{\infty}$]
    \label{lem:M_i_in_M_inf}
    \lean{Construction.M_I_in_M_inf,Construction.M_I_in_M_inf_elm}
    \leanok
    \uses{def:set_of_constructable_points}
    The set $\mathcal{M}_i$ is in $\mathcal{M}_{\infty}$, i.e. $$\mathcal{M}_i \subseteq \mathcal{M}_{\infty}.$$
\end{lemma}
\begin{proof}
    \leanok
    Follows from the definition of $\mathcal{M}_{\infty}$.
\end{proof}

\begin{lemma}[$\mathcal{M}$ in $\mathcal{M}_{\infty}$]
    \label{lem:M_in_M_inf}
    \lean{Construction.M_M_inf}
    \leanok
    \uses{def:set_of_constructable_points}
    The set $\mathcal{M}$ is in $\mathcal{M}_{\infty}$.
\end{lemma}
\begin{proof}
    \leanok
    \uses{lem:M_i_in_M_inf, lem:M_in_Mi}
    Combining $\mathcal{M} \subseteq \mathcal{M}_i$ \ref{lem:M_in_Mi} and $\mathcal{M}_i \subseteq \mathcal{M}_{\infty}$ \ref{lem:M_i_in_M_inf} we get the result.
\end{proof}

\begin{lemma}[$\mathcal{M}_{\infty}$ iff $\mathcal{M}_i$]
    \label{lem:M_inf_iff_M_i}
    \lean{Construction.M_inf_in_M_I, Construction.M_inf_in_M_I'}
    \leanok
    \uses{def:set_of_constructable_points}
    A point $z$ is in $\mathcal{M}_{\infty}$ if and only if $z$ is in $\mathcal{M}_i$ for some $i$.
\end{lemma}
\begin{proof}
    \leanok
    Follows from the definition of $\mathcal{M}_{\infty}$.
\end{proof}

\begin{corollary}[$L(\mathcal{M_{\infty}})$ iff $L(\mathcal{M_i})$]
    \label{cor:L_M_inf_iff_L_M_i}
    \lean{Construction.L_M_inf_iff_M_M}
    \leanok
    \uses{def:set_of_lines, def:set_of_constructable_points}
    A line $l$ is in $\mathcal{L(M_{\infty})}$ if and only if $l$ is in $\mathcal{L(M_i)}$ for some $i$.
\end{corollary}
\begin{proof}
    \leanok
    \uses{lem:M_inf_iff_M_i}
    Follows from \ref{lem:M_inf_iff_M_i} and the fact that every line in $\mathcal{L(M_{\infty})}$ is defined by two points in $\mathcal{M_{\infty}}$.
\end{proof}

\begin{corollary}[$C(\mathcal{M_{\infty}})$ iff $C(\mathcal{M_i})$]
    \label{cor:C_M_inf_iff_C_M_i}
    \lean{Construction.C_M_inf_iff_M_M}
    \leanok
    \uses{def:set_of_circles, def:set_of_constructable_points}
    A circle $c$ is in $\mathcal{C(M_{\infty})}$ if and only if $c$ is in $\mathcal{C(M_i)}$ for some $i$.
\end{corollary}
\begin{proof}
    \leanok
    \uses{lem:M_inf_iff_M_i}
    Follows from \ref{lem:M_inf_iff_M_i} and the fact that every circle in $\mathcal{C(M_{\infty})}$ is defined by three points in $\mathcal{M_{\infty}}$.
\end{proof}

To construct new points in the next section\ref{basic_constructions} we still need to show that every intersection of lines and circles is in $\mathcal{M}_{\infty}$.
\begin{lemma}[Intersection of lines in $\mathcal{M}_{\infty}$]
    \label{lem:ill_M_inf}
    \lean{Construction.ill_M_inf}
    \leanok
    \uses{def:set_of_constructable_points, def:set_of_lines}
    For two lines $l_1, l_2 \in \mathcal{L(M_{\infty})}$ is $l_1 \cap l_2 \in \mathcal{M}_{\infty}$.
\end{lemma}
\begin{proof}
    \leanok
    \uses{cor:L_M_inf_iff_L_M_i, lem:M_i_monoton, lem:M_i_in_M_inf}
    By corollary \ref{cor:L_M_inf_iff_L_M_i} and the monotonity of $\mathcal{M}_i$ \ref{lem:M_i_monoton} we get that there exists an $n$ such that $l_1, l_2 \in \mathcal{L(M_n)}$.
    Therfore $l_1 \cap l_2 \in \mathcal{M}_{n+1} \stackrel{\ref{lem:M_i_in_M_inf}}{\subseteq} \mathcal{M}_{\infty}$.
\end{proof}

\begin{remark}%TODO: check if remark is needed
    To formalise the lemmas \ref{lem:M_inf_iff_M_i}, \ref{cor:L_M_inf_iff_L_M_i} and \ref{cor:C_M_inf_iff_C_M_i} one should use filters in Lean.
\end{remark}

\begin{lemma}[Intersection of line and circle in $\mathcal{M}_{\infty}$]
    \label{lem:ilc_M_inf}
    \lean{Construction.ilc_M_inf}
    \leanok
    \uses{def:set_of_constructable_points, def:set_of_lines, def:set_of_circles}
    For a line $l \in \mathcal{L(M_{\infty})}$ and a circle $c \in \mathcal{C(M_{\infty})}$ is $l \cap c \in \mathcal{M}_{\infty}$.
\end{lemma}
\begin{proof}
    \leanok
    \uses{cor:L_M_inf_iff_L_M_i, cor:C_M_inf_iff_C_M_i, lem:M_i_monoton, lem:M_i_in_M_inf}
    By corollary \ref{cor:L_M_inf_iff_L_M_i} and \ref{cor:C_M_inf_iff_C_M_i} we get that there exists an $n$ such that $l \in \mathcal{L(M_n)}$ and $c \in \mathcal{C(M_n)}$.
    Therfore $l \cap c \in \mathcal{M}_{n+1} \stackrel{\text{\ref{lem:M_i_in_M_inf}}}{\subseteq} \mathcal{M}_{\infty}$.
\end{proof}

\begin{lemma}[Intersection of circles in $\mathcal{M}_{\infty}$]
    \label{lem:icc_M_inf}
    \lean{Construction.icc_M_inf}
    \leanok
    \uses{def:set_of_constructable_points, def:set_of_circles}
    For two circles $c_1, c_2 \in \mathcal{C(M_{\infty})}$ it follows that $c_1 \cap c_2 \in \mathcal{M}_{\infty}$.
\end{lemma}
\begin{proof}
    \leanok
    \uses{cor:C_M_inf_iff_C_M_i, lem:M_i_monoton, lem:M_i_in_M_inf}
    By corollary \ref{cor:C_M_inf_iff_C_M_i} we get that there exists an $n$ such that $c_1, c_2 \in \mathcal{C(M_n)}$.
    Therfore $c_1 \cap c_2 \in \mathcal{M}_{n+1} \stackrel{\ref{lem:M_i_in_M_inf}}{\subseteq} \mathcal{M}_{\infty}$.
\end{proof}


\section{Basic constructions}\label{basic_constructions}
It is now possible to construct fundamental points in $M_{\infty}$ using a compass and ruler, which can subsequently be employed to create a field structure and special properties on $M_{\infty}$. 
Consequently, the following constructions are based on the assumption that $M\subseteq \mathbb{C}$ with $0,1 \in M$.
%2.3.1
\begin{lemma}[Negative complex numbers]
    \label{lem:construction_neg}
    \lean{Construction.z_neg_M_inf}
    \leanok
    \uses{def:set_of_constructable_points}
    For $z \in M_{\infty}$ it follows that $-z \in M_{\infty}$.
\end{lemma}
This construction is taken from \cite{JAN_SCHRÖER:2023}.\\
To get the point $-z$ we can use the second intersection of the line through $0$ and $z$ with the circle with center $0$ and radius $\|z\|$.\ref{Fig.1}

\begin{proof}
    \leanok
    \uses{def:line, def:circle, lem:ilc_M_inf, lem:M_in_M_inf}
    Define $l = \{0,z\}$ and $c = \{0,\dist (0, z)\}$.\\
    By assumption $0, z \in M_{\infty}$, so $l \in \mathcal{L(M_{\infty})}$ and $c \in \mathcal{C(M_{\infty})}$.\\
    \underline{Claim 1}: $-z$ is in $l$.
    \begin{proof}
        By the definition of $l:=\{\lambda 0+(1-\lambda)z\mid\lambda\in\mathbb{R}\}$ with $\lambda = 2$ we get $2 \cdot 0 + (1-2)z = -z$.
    \end{proof}
    \underline{Claim 2}: $-z$ is in $c$.
    \begin{proof}
        Unfolding the definition of $c:=\{x\in\mathbb{C} \mid\|x-0\|=\dist (0\ z)\}$.
        By the definition of the distance we get $\|0-(-z)\| = \dist (0\ z)$.
    \end{proof}
    By claim 1 and 2 we get that $-z \in l \cap c$. Futhermore  $-z $ is in $M_{\infty}$, after lemma \ref{lem:ilc_M_inf}.
\end{proof}

\begin{figure}[h!]
    \centering
    \begin{tikzpicture}
        \draw[->] (-1,0) -- (3,0) node[right] {$\Re$};
        \draw[->] (0,-1) -- (0,3) node[above] {$\Im$};
        \draw (0,0) -- (2,2) node[right] {$z$};
        \draw (0,0) -- (-2,-2) node[left] {$-z$};
        \draw (0,0) circle (2.8);
    \end{tikzpicture}
    \caption{Construction of $-z$}
    \label{Fig.1}
\end{figure}

%2.3.2
\begin{lemma}[Addition of complex numbers]
    \label{lem:construction_add}
    \lean{Construction.add_M_Inf}
    \leanok
    \uses{def:set_of_constructable_points}
    For $z_1, z_2 \in M_{\infty}$ it follows that $z_1 + z_2 \in M_{\infty}$.
\end{lemma}
This construction is taken from \cite{JAN_SCHRÖER:2023}.\\
One can construct the point $z_1 + z_2$ by drawing a circle with center $z_1$ and radius $\|z_2\|$ and a circle with center $z_2$ and radius $\|z_1\|$ and taking the intersection of the two circles.\ref{Fig.2}
\begin{proof}
    \leanok
    \uses{def:line, def:circle,  lem:icc_M_inf, lem:M_in_M_inf}
    First we have the case that $z_1 \ne z_2$.\\
    Define $c_1 = \{z_1,\dist (0\ z_2)\}$ and $c_2 = \{z_2,\dist (0\ z_1)\}$.\\
    By assumption $0, z_1, z_2 \in M_{\infty}$, so $c_1, c_2 \in \mathcal{C(M_{\infty})}$ and since $z_1 \ne z_2$ we get $c_1 \ne c_2$.\\
    \underline{Claim 1}: $z_1 + z_2$ is in $c_1$.
    \begin{proof}
        By the definition of $c_1:=\{x\in\mathbb{C} \mid\|x-z_1\|=\dist (0\ z_2)\}$
        and of the distance we get $\|z_1-(z_1+z_2)\| = \dist (0\ z_2)$.
    \end{proof}
    \underline{Claim 2}: $z_1 + z_2$ is in $c_2$.
    \begin{proof}
        Using that  $c_2$ ia definded as $c_2:=\{x\in\mathbb{C} \mid\|x-z_2\|=\dist (0\ z_1)\}$ and the definition of the distance we get $\|z_2-(z_1+z_2)\| = \dist (0\ z_1)$.
    \end{proof}
    By claim 1 and 2 we get that $z_1 + z_2 \in c_1 \cap c_2$. Futhermore  $z_1 + z_2 $ is in $M_{\infty}$, after lemma \ref{lem:icc_M_inf}.
    
    If we have $z_1 = z_2$ we can define $c_1 = \{z_1,\dist (0\ z_1)\}$ and $l = \{0,z_1\}$ and get that $z_1 + z_2 = z_1 + z_1 \in c_1 \cap l$.

\end{proof}
% \begin{figure}[h]
%     \centering
%     \begin{tikzpicture}
%         \draw[->] (-1,0) -- (3,0) node[right] {$\Re$};
%         \draw[->] (0,-1) -- (0,3) node[above] {$\Im$};
%         \draw (0,0) -- (1,2) node[above] {$z_1$}[red];
%         \draw (0,0) -- (2,1) node[right] {$z_2$}[blue];
%         \draw (1,2) circle (2.2360679775)[red];
%         \draw (2,1) circle (2.2360679775)[blue];
%         \draw[dashed,->] (1,2) -- (3,3)[blue];
%         \draw[dashed,->] (2,1) -- (3,3)[red];
%         \draw (0,0) -- (3,3) node[right] {$z_1 + z_2$};
%     \end{tikzpicture}
%     \caption{Construction of $z_1 + z_2$}
%     \label{Fig.2}
% \end{figure}

\begin{figure}[h]
    \centering
    \begin{tikzpicture}
        \draw[->] (-1,0) -- (3,0) node[right] {$\Re$};
        \draw[->] (0,-1) -- (0,3) node[above] {$\Im$};
        \draw (0,0) -- (1,2) node[above] {$z_1$};
        \draw (0,0) -- (2,1) node[right] {$z_2$};
        \draw (1,2) circle (2.2360679775);
        \draw (2,1) circle (2.2360679775);
        \draw[dashed,->] (1,2) -- (3,3);
        \draw[dashed,->] (2,1) -- (3,3);
        \draw (0,0) -- (3,3) node[right] {$z_1 + z_2$};
    \end{tikzpicture}
    \caption{Construction of $z_1 + z_2$}
    \label{Fig.2}
\end{figure}

\begin{corollary}[Subtraction of complex numbers]
    \label{cor:construction_sub}
    \lean{Construction.sub_M_Inf}
    \leanok
    \uses{def:set_of_constructable_points}
    For $z_1, z_2 \in M_{\infty}$ it follows that $z_1 - z_2 \in M_{\infty}$.
\end{corollary}
\begin{proof}
    \leanok
    \uses{lem:construction_neg, lem:construction_add}
    By lemma \ref{lem:construction_neg} and \ref{lem:construction_add} we get $z_1 - z_2 = z_1 + (-z_2) \in M_{\infty}$.
\end{proof}

\begin{corollary}[Construction of parallel lines]
    \label{lem:construction_parallel_lines}
    \lean{Construction.parallel_lines_M_inf}
    \leanok
    \uses{def:set_of_constructable_points, def:set_of_lines, def:parallel}
    For $l \in \mathcal{L(M_{\infty})}$ and $z \in M_{\infty}$ it follows that $\exists l' \in \mathcal{L(M_{\infty})}$ with  $z\in l'$  and  $l\|l'$.
\end{corollary}
\begin{proof}
    \leanok
    \uses{def:line, def:set_of_lines, cor:construction_sub,lem:parallel_by_basis} 
    Let $l$ be a line through $x,y \in \mathcal{M_{\infty}}$ and $z \in \mathcal{M_{\infty}}$.\\
    After corollary \ref{cor:construction_sub} we get $z - x \in \mathcal{M_{\infty}}$ and therfore $z-x + y \in \mathcal{M_{\infty}}$\ref{lem:construction_add}.\\
    Define $l' = \{z, z-x+y\}$, then $l' \in \mathcal{L(M_{\infty})}$ and since a line is defined by two points and we moved them the same distance ($z-x$) $l'$ is parallel to $l$.\ref{lem:parallel_by_basis}
\end{proof}
%2.3.3
%TODO: Comper to mirrowed version
\begin{lemma}[Complex conjugation]
    \label{lem:construction_conj}
    \lean{Construction.conj_M_Inf}
    \leanok
    \uses{def:set_of_constructable_points}
    For $z \in M_{\infty}$ it follows that $\overline{z} \in M_{\infty}$.
\end{lemma}
This construction is taken from \cite{JAN_SCHRÖER:2023}.\\
Draw two circles one with center $0$ and radius $\|z\|$ and a second with center $1$ and radius $\|z-1\|$ and take the intersection of the two circles.\ref{Fig.3}
\begin{proof}
    \leanok
    \uses{def:line, def:circle, lem:icc_M_inf, lem:M_in_M_inf}
    Define $c_1 = \{0,\dist (0, z)\}$ and $c_2 = \{1,\dist (1, z)\}$.\\
    By assumption $0, 1, z \in M_{\infty}$, so $c_1, c_2 \in \mathcal{C(M_{\infty})}$.\\
    \underline{Claim 1}: $\overline{z}$ is in $c_1$.
    \begin{proof}
        The definition of $c_1$ can be written in the following form: $$\{x\in\mathbb{C} \mid \|x-0\| = \dist (0\ z)\}.$$ 
        Using the definition of the distance, we obtain the following result: $$\|0 - \overline{z}\| = \|\overline{z}\| = \|z\| = \dist (0\ z).$$
    \end{proof}
    \underline{Claim 2}: $\overline{z}$ is in $c_2$.
    \begin{proof}
        From the definition of $c_2:=\{x\in\mathbb{C} \mid\|x-1\|=\dist (1\ z)\}$ and the definition of the distance, we can derive the following:
         $$\|1-\overline{z}\| = \|\overline{z}-1\| = \sqrt{(\Re(\overline{z})-1)^2 + \Im(\overline{z})^2} = \sqrt{(\Re(z)-1)^2 + \Im(z)^2} = \|z-1\|=\dist (1\ z).$$
    \end{proof}
    By claim 1 and 2 we get that $\overline{z} \in c_1 \cap c_2$. Futhermore  $\overline{z} $ is in $M_{\infty}$, after lemma \ref{lem:icc_M_inf}.
\end{proof}
\begin{figure}[h]
    \centering
    \begin{tikzpicture}
        \draw[->] (-3,0) -- (3,0) node[right] {$\Re$};
        \draw[->] (0,-2.5) -- (0,2.5) node[above] {$\Im$};
        \draw[dashed] (0,0) -- (2,1) node[right] {$z$};
        \draw[dashed] (0,0) -- (2,-1) node[right] {$\overline{z}$};
        \draw (0,0) circle (2.2360679775);
        \draw (1,0) circle (1.4142135624);
    \end{tikzpicture}
    \caption{Construction of $\overline{z}$}
    \label{Fig.3}
\end{figure}


\begin{lemma}[Construction $\imath r\in\mathcal{M}_{\infty}$]
    \label{lem:construction_imath_r}
    \lean{Construction.ir_M_inf}
    \leanok
    \uses{def:set_of_constructable_points}
    For $r \in \mathbb{R}\cap M_{\infty}$ it follows that $\imath\cdot r \in M_{\infty}$.
\end{lemma}
The first step is to construct the imaginary axis, which is achieved by drawing a line, designated as $l$, that passes through two circles with centres at $-1$ and $1$ and radii of $2$. 
The second step is to construct the point, designated as $\imath r$, which is achieved by drawing a circle with a centre at $0$ and a radius of $r$, and taking the intersection with the imaginary axis. (Fig\ref{Fig.4})
\begin{proof}
    \leanok
    \uses{def:line, def:circle, lem:construction_neg, lem:icc_M_inf, lem:ilc_M_inf, lem:M_in_M_inf}
    First we construct the imaginary axis. Define $c_1 = \{-1,2\}$ and $c_2 = \{1,2\}$.\\
    \underline{Claim 1}: $c_1, c_2 \in \mathcal{C(M_{\infty})}$
    \begin{proof}
        By assumption and lemma \ref{lem:construction_neg} we get $-1, 1 \in M_{\infty}$. Using $\dist(-1 1)=2$ we get $c_1, c_2 \in \mathcal{C(M_{\infty})}$, by the definition of the circle.
    \end{proof}
    \underline{Claim 2}: $\imath\sqrt{3}$ and $-\imath\sqrt{3}$ are in $c_1 \cap c_2$.
    \begin{proof}
        Unfolding the definition of $c_1:=\{x\in\mathbb{C} \mid\|x+1\|=2\}$ and $c_2:=\{x\in\mathbb{C} \mid\|x-1\|=2\}$.
        By the definition of the distance we get $\|-\imath\sqrt{3}+1\| = \sqrt{1 + 3} = 2$ and $\|\imath\sqrt{3}-1\| = \sqrt{1 + 3} = 2$.
    \end{proof}
    Now we define $l = \{\imath\sqrt{3}, -\imath\sqrt{3}\}$.\\
    \underline{Claim 3}: $l \in \mathcal{L(M_{\infty})}$
    \begin{proof}
        By claim 2 and \ref{lem:icc_M_inf} we get $\imath\sqrt{3}, -\imath\sqrt{3} \in M_{\infty}$, so $l \in \mathcal{L(M_{\infty})}$.
    \end{proof}
    To get $\imath r$ we define $c = \{0,|r|\}$.\\
    \underline{Claim 4}: $\imath r$ is in $c \cap l$.
    \begin{proof}
        It is clear that $\imath r \in c$. Now using the definition of $l$ and $\lambda = \frac{r}{2\sqrt{3}+\frac{1}{2}}$ we get $(\frac{r}{2\sqrt{3}+\frac{r}{2}})\imath\sqrt{3} + (1-\frac{r}{2\sqrt{3}+\frac{r}{2}})(-\imath\sqrt{3}) = \imath r$. 
    \end{proof}
    Therfore $\imath r \in M_{\infty}$ after lemma \ref{lem:ilc_M_inf}.
\end{proof}
\begin{figure}[h]
    \centering
    \begin{tikzpicture}  
        \draw[->] (-3,0) -- (3,0) node[right] {$\Re$};
        \draw[->] (0,-2.5) -- (0,2.5) node[above] {$\Im$};
        \coordinate[label=-45:$1$] (one) at (1,0);
        \coordinate[label=-135:$-1$] (mone) at (-1,0);
        \fill[black] (one) circle (2pt);
        \fill[black] (mone) circle (2pt);
        \draw[dashed] (one) [carc=110:250:2];
        \draw[dashed] (mone) [carc=70:-70:2];
        \coordinate[label=45:$r$] (r) at (2,0);
        \coordinate[label=135:$\imath r$] (ir) at (0,2);
        \fill[black] (r) circle (2pt);
        \fill[black] (ir) circle (2pt);
        \draw (0,0) [carc=-10:100:2];
    \end{tikzpicture}
    \caption{Construction of $\imath r$}
    \label{Fig.4}
\end{figure}

\begin{corollary}[Construction of $\imath$]
    \label{cor:construction_imath}
    \lean{Construction.imath_M_inf}
    \leanok
    \uses{def:set_of_constructable_points}
    $\imath \in M_{\infty}$.
\end{corollary}
\begin{proof}
    \leanok
    \uses{lem:construction_imath_r, lem:M_in_M_inf}
    By lemma \ref{lem:construction_imath_r} with $r = 1$ and the definition of $\mathcal{M}$ we get $\imath \in M_{\infty}$.
\end{proof}

\begin{lemma}[Construction of real part]
    \label{lem:construction_real}
    \lean{Construction.real_in_M_inf}
    \leanok
    \uses{def:set_of_constructable_points}
    For $z \in M_{\infty}$ it follows that $z.re \in M_{\infty}$.
\end{lemma}
To get the point $z.re$ we draw a line through $z$ and $\overline{z}$. Then $z.re$ is the intersection with the real line, defined by $0$ and $1$.\ref{Fig z.re}
\begin{proof}
    \leanok
    \uses{def:line, } %TODO c_in_C_M, ilc_M_inf, M_M_inf
    Without loss of generality we can assume that $z \in \C\setminus\R$.\\
    Define the lines $l = \{z, \overline{z}\}$ and $l_{\Re} = \{1, 0\}$.\\
    By using Lemma \ref{lem:construction_conj}, they are in $\mathcal{L(M_{\infty})}$.\\
    To show that $z.re \in l \cap l_{\Re}$ we use $t:= 1/2$ for $l$ and $t:= z.re$ for $l_{\Re}$.
\end{proof}
\begin{figure}
    \centering
    \begin{tikzpicture}
        \draw[->] (-0.5,0) -- (5,0) node[right] {$\Re$};
        \draw[->] (0,-3) -- (0,3) node[above] {$\Im$};
        \coordinate[label=135:$z$] (z) at (2,2);
        \coordinate[label=-45:$z.re$] (zre) at (2,0);
        \coordinate[label=135:$\overline{z}$] (oz) at (2,-2);
        \coordinate[label=-135:$0$] (zero) at (0,0);
        \coordinate[label=-90:$1$] (one) at (1,0);
        \fill[black] (z) circle (2pt);
        \fill[black] (zre) circle (2pt);
        \fill[black] (oz) circle (2pt);
        \fill[black] (one) circle (1pt);
        \fill[black] (zero) circle (1pt);
        \draw (z) -- (oz);
    \end{tikzpicture}
    \caption{Construction of $z.re$}
    \label{Fig z.re}
\end{figure}

\begin{lemma}[Construction of imaginary part]
    \label{lem:construction_imag}
    \lean{Construction.im_in_M_inf,Construction.i_z_imp_z_in_M_inf}
    \leanok
    \uses{def:set_of_constructable_points}
    For $z \in M_{\infty}$ it follows that $z.im \in M_{\infty}$.
\end{lemma}
To get the point $z.im$ we draw a line through $z$ and $z-1$. Now we get $\imath \cdot z.im$ by taking the intersection with the imaginary line, defined by $0$ and $\imath$. To get $z.im$ draw a circle trough $\imath \cdot z.im$ and take the intersection with the real line.\ref{Fig z.im}
\begin{proof}
    \leanok
    \proves{lem:construction_imag}
    \uses{def:line, def:circle, c_in_C_M, ilc_M_inf, M_M_inf}%TODO 
    Define the lines $l = \{z, z-1\}$ and $l_{\Im} = \{\imath, 0\}$.\\
    Now it is analog to the proof of lemma \ref{lem:construction_real}, that $z.im \cdot \imath  \in l \cap l_{\Im}$. 
    %TODO: Remove +0 if issue is fixed
    Now we can project $z.im \cdot \imath + 0$ on the real axis by drawing a circle with center $0$ and radius $\|z.im\|$ and taking the intersection with the real axis.
\end{proof}
\begin{figure}
    \centering
    \begin{tikzpicture}
        \draw[->] (-1,0) -- (4,0) node[right] {$\Re$};
        \draw[->] (0,-1) -- (0,4) node[above] {$\Im$};
        \coordinate[label=90:$z$] (z) at (2,3);
        \coordinate[label=-45:$z.im$] (zim) at (3,0);
        \coordinate[label=90:$z-1$] (z1) at (1,3);
        \coordinate[label=-135:$\imath\cdot z.im$] (izim) at (0,3);
        \coordinate[label=-135:$0$] (zero) at (0,0);
        \coordinate[label=180:$\imath$] (i) at (0,1);
        \coordinate[label=-90:$1$] (one) at (1,0);
        \fill[black] (z) circle (2pt);
        \fill[black] (izim) circle (2pt);
        \fill[black] (zim) circle (2pt);
        \fill[black] (z1) circle (2pt);
        \fill[black] (i) circle (1pt);
        \fill[black] (zero) circle (1pt);
        \fill[black] (one) circle (1pt);
        \draw (-0.5,3) -- (3.5,3);
        \draw (0,0) [carc=-30:120:3];
    \end{tikzpicture}
    \caption{Construction of $z.im$}
    \label{Fig z.im}
\end{figure}

\begin{corollary}[$z \in M_{\infty} \Rightarrow z.re, z.im \in M_{\infty}$]
    \label{cor:z_iff_re_im_M_inf}
    \lean{Construction.z_iff_re_im_M_inf}
    \leanok
    \uses{def:set_of_constructable_points}
    For $z \in \C$. $z$ is in $M_{\infty}$ if and only if $z.re, z.im \in M_{\infty}$.
\end{corollary}
\begin{proof}
    \leanok
    \uses{lem:construction_real, lem:construction_imag, lem:construction_add, lem:construction_imath_r}
    "$\Rightarrow$:" If $z \in M_{\infty}$ then $z.re, z.im \in M_{\infty}$ after lemma \ref{lem:construction_real} and \ref{lem:construction_imag}.\\
    "$\Leftarrow$:" If $z.re, z.im \in M_{\infty}$ then $z.re + \imath z.im \in M_{\infty}$ after lemma \ref{lem:construction_add} and lemma \ref{lem:construction_imath_r}.
\end{proof}

\begin{lemma}[Multiplication of positve real numbers]
    \label{lem:construction_mul}
    \lean{Construction.ab_in_M_inf}
    \leanok
    \uses{def:set_of_constructable_points}
    For $a, b \in M_{\infty}\cap\R$ it follows that $a \cdot b \in M_{\infty}$.
\end{lemma}
This construction is taken from \cite{cox2012galois}.\\
To get the point $a\cdot b$ we draw a line trough $a$ and $\imath$ and a parallel line through $\imath b$. The intersection of the second line with the real axis is $a\cdot b$.\ref{Fig.5} 
\begin{proof}
    \leanok
    \uses{def:line, def:set_of_lines, lem:ill_M_inf, lem:construction_imath_r, cor:construction_imath, lem:construction_add, cor:construction_sub}
    Define the three lines $l = \{a+\imath b - \imath, \imath b\}$ and $l_{\Re} = \{1,0\}$.\\
    \underline{Claim 1}: $l \in \mathcal{L(M_{\infty})}$
    \begin{proof}
        By assumption $b \in M_{\infty}$ so after Lemma \ref{lem:construction_imath_r} $\imath b \in M_{\infty}$. Futhermore $l_2 \in \mathcal{L(M_{\infty})}$ after Claim 1 and Lemma \ref{lem:construction_parallel_lines}.
    \end{proof}
    \underline{Claim 2}: $l_{\Re} \in \mathcal{L(M_{\infty})}$
    \begin{proof}
        By assumption $0, 1 \in M_{\infty}$, so $l_{\Re} \in \mathcal{L(M_{\infty})}$.
    \end{proof}
    To show that $a\cdot b \in M_{\infty}$ we need to show that $l \cap l_{\Re} \in M_{\infty}$\ref{lem:ill_M_inf}. That $ab \in l_{\Re} \cap l$ is clear after definition $t=ab$. For $ab \in l$ we use $t=b$ and get $t(a+\imath b - \imath) + (1-t)\imath b  \stackrel{t:=b}{=}ba + \imath b^2 -\imath b + \imath b -\imath b^2 = a\cdot b$.
\end{proof}
\begin{remark}
    This construction uses parallel lines, but it is not needed for the proof of the lemma.
\end{remark}
\begin{figure}[!h]
    \centering
    \begin{tikzpicture}
        \draw[->] (-0.5,0) -- (5,0) node[right] {$\Re$};
        \draw[->] (0,-0.5) -- (0,3) node[above] {$\Im$};
        \coordinate[label=135:$\imath$] (i) at (0,1);
        \coordinate[label=-90:$a$] (a) at (2,0);
        \coordinate[label=135:$\imath b$] (ib) at (0,2);
        \coordinate[label=-90:$ab$] (ab) at (4,0);
        \fill[black] (i) circle (2pt);
        \fill[black] (a) circle (2pt);
        \fill[black] (ib) circle (2pt);
        \fill[black] (ab) circle (2pt);
        \draw (a) -- (i);
        \draw (ib) -- (ab);
    \end{tikzpicture}
    \caption{Construction of $z_1 \cdot z_2$}
    \label{Fig.5}
\end{figure}

\begin{corollary}[Multiplication of complex numbers]
    \label{cor:construction_mul_complex}
    \lean{Construction.mul_M_inf} 
    \leanok
    \uses{def:set_of_constructable_points}
    For $z_1, z_2 \in M_{\infty}$ it follows that $z_1 \cdot z_2$ in $M_{\infty}$.
\end{corollary}
\begin{proof}
    \leanok
    \uses{lem:construction_mul, lem:construction_add, cor:construction_sub, lem:construction_real, lem:construction_imag}
    Let $z_1 = a + \imath b$ and $z_2 = c + \imath d$. Then $$z_1 \cdot z_2 = (a + \imath b) \cdot (c + \imath d) = (a \cdot c - b \cdot d) + \imath (a \cdot d + b \cdot c).$$
    By combining the Lemmas \ref{lem:construction_mul}, \ref{lem:construction_add}, \ref{cor:construction_sub}, \ref{lem:construction_real} and \ref{lem:construction_imag} we get that $z_1 \cdot z_2 \in M_{\infty}$.
\end{proof}

\begin{lemma}[Invers of a pos real number]
    \label{lem:construction_inv}
    \lean{Construction.ainv_in_M_inf}
    \leanok
    \uses{def:set_of_constructable_points}
    If $a \in M_{\infty}\cap\R$, then $a^{-1}$ is in  $M_{\infty}$.
\end{lemma}

This can be constructed analog to the multiplication of positve real numbers. 
Using the fact that $a\cdot a^{-1} = 1$.
Draw a line through $1$ and $\imath a$ and a parallel line through $\imath$. 
The intersection of the second line with the real axis is $a^{-1}$. (Fig. \ref{Fig.6})
\begin{proof}
    \leanok
    \uses{def:line, def:set_of_lines, lem:ill_M_inf, lem:construction_imath_r, cor:construction_imath, lem:construction_add, cor:construction_sub}
    The proof is analogous to that of Lemma \ref{lem:construction_mul}. It requires only two lines: $l = \{1 - \imath z + \imath, \imath\}$ and $l_{\Re} = \{1,0\}$.\\
    With  out loss of generality we can assume that $a \ne 0$.\\
    The fact that they are in $\mathcal{L(M_{\infty})}$ follows analog to the proof of Lemma \ref{lem:construction_mul}.\\ 
    Thus we have just to show that $z^{-1} \in l$, i.e. $\exists t: t  (1 - \imath a + \imath) + (1 - t)  I = a^{-1}$ $$t  (1 - \imath a + \imath) + (1 - t)  \imath \stackrel{t:=a^{-1}}{=}  a^{-1} - a^{-1} \imath a + a^{-1}\imath + \imath - a^{-1}\imath = a^{-1}.$$
    The rest follows analog.
\end{proof}
\begin{figure}
    \centering
    \begin{tikzpicture}
        \draw[->] (-1,0) -- (2,0) node[right] {$\Re$};
        \draw[->] (0,-1) -- (0,3) node[above] {$\Im$};
        \coordinate[label=135:$\imath$] (i) at (0,1);
        \coordinate[label=-90:$1$] (1) at (1,0);
        \coordinate[label=135:$\imath a$] (ia) at (0,2);
        \coordinate[label=-90:$a^{-1}$] (ainv) at (0.5,0);
        \fill[black] (i) circle (2pt);
        \fill[black] (1) circle (2pt);
        \fill[black] (ia) circle (2pt);
        \fill[black] (ainv) circle (2pt);
        \draw (1) -- (ia);
        \draw (i) -- (ainv);
    \end{tikzpicture}
    \caption{Construction of $z^{-1}$}
    \label{Fig.6}
\end{figure}

\begin{corollary}[Inverse of a complex number]
    \label{cor:inv_M_inf}
    \lean{Construction.inv_M_inf, Construction.z_inv_eq}
    \leanok
    \uses{def:set_of_constructable_points}
    If $z \in M_{\infty}$, then $z^{-1}$ is in  $M_{\infty}$.
\end{corollary}
\begin{proof}
    \leanok
    \uses{lem:construction_inv, lem:construction_mul, cor:construction_sub, lem:construction_real, lem:construction_imag, lem:construction_add}
    For $z \in M_{\infty}$ we can write $z = a + \imath b$ with $a, b \in \R$. Then
    $$z^{-1} = \frac{1}{z} = \frac{\overline{z}}{z\overline{z}} = \frac{a - \imath b}{a^2+b^2}= (a - \imath b) \cdot (aa+bb)^{-1}.$$
    It is now possible to combine the lemmas for addition \ref{lem:construction_add}, subtraction \ref{cor:construction_sub}, multiplication \ref{cor:construction_mul_complex} 
    and the corollary for the inverse of a positive real number \ref{lem:construction_inv} with the part concerning the existence of real and imaginary components \ref{cor:z_iff_re_im_M_inf} in order to conclude that $z^{-1} \in M_{\infty}$.
\end{proof}


\begin{lemma}[Angle in $M_{\infty}$]
    \label{lem:angle_M_inf}
    \uses{def:set_of_constructable_points}
    \lean{Construction.angle_M_inf}
    \leanok
    If $0 \ne z = r\cdot e^{\imath\alpha} \in M_{\infty}$, then $e^{\imath\alpha} \in  M_{\infty}$.
\end{lemma}
For the construction we draw a line through $0$ and $z$ and take the intersection with the unit circle. \ref{Fig.angle}

\begin{proof}
    \leanok
    \uses{lem:icc_M_inf}
    Let $l$ be a line through $0$ and $z$ and $c$ be the unit circle. Then $l \cap c = \{e^{\imath\alpha}\}$.
    The rest follows from the construction of the intersection of a line and a circle\ref{lem:icc_M_inf}.
\end{proof}

\begin{figure}[h]
    \centering
    \begin{tikzpicture}
        \draw[->] (-0.5,0) -- (4,0) node[right] {$\Re$};
        \draw[->] (0,-0.5) -- (0,3) node[above] {$\Im$};
        \coordinate[label=-135:$0$] (zero) at (0,0);
        \coordinate[label=-135:$1$] (one) at (1,0);
        \coordinate[label=135:$z$] (z) at (3,2);
        \coordinate[label=90:$e^{\imath\alpha}$] (ea) at (0.8321,0.5547);
        \fill[black] (zero) circle (1pt);
        \fill[black] (one) circle (2pt);
        \fill[black] (z) circle (2pt);
        \fill[black] (ea) circle (2pt);
        \draw (zero) -- (z);
        \draw (zero) [carc=-20:110:1];
    \end{tikzpicture}
    \caption{Construction of $e^{\imath\alpha}$}
    \label{Fig.angle}
\end{figure}

\begin{corollary}[Midpoint in $M_{\infty}$]
    \label{cor:Construction.midpoiont}
    \uses{def:set_of_constructable_points}
    \lean{Construction.midpoiont}
    \leanok
    If $z_1, z_2 \in M_{\infty}$, then $\frac{z_1 + z_2}{2} \in  M_{\infty}$.
    
\end{corollary}

\begin{proof}
    \uses{lem:construction_add, lem:construction_mul, lem:construction_inv}
    \leanok
    Combining the lemmas for addition\ref{lem:construction_add}, multiplication\ref{cor:construction_mul_complex} and the invers of a complex number\ref{cor:inv_M_inf} we get that $\frac{z_1 + z_2}{1+1} \in M_{\infty}$.
    
\end{proof}

\begin{lemma}[Halving of an angle]
    \label{lem:Construction.angle_half_M_inf}
    \uses{def:set_of_constructable_points}
    \lean{Construction.angle_half_M_inf}
    \leanok
    For $\alpha \in [0,2\pi)$, if $e^{\imath\alpha} \in M_{\infty}$, then $e^{\imath\frac{\alpha}{2}} \in  M_{\infty}$.
\end{lemma}
\begin{proof}
    \leanok
    \uses{lem:angle_M_inf, lem:icc_M_inf}
    For $\alpha \ne 0 \ne \pi$ we take the intersection of the unit circle with the line through $0$ and and the mipiont of $e^{\imath\alpha}$ and $1$. 
    %TODO: Remove +0 if issue is fixed
    For $\alpha = 0$ we get $e^{\imath\frac{\alpha}{2}} = 1$ and for $\alpha = \pi$ we get $e^{\imath\frac{\alpha}{2}} =\imath + 0$.
\end{proof}

\begin{lemma}[Construction of radius]
    \label{lem:construction_radius}
    \uses{def:set_of_constructable_points}
    \lean{Construction.abs_M_Inf}
    \leanok
    If $z = r\cdot e^{\imath\alpha} \in M_{\infty}$, then $r \in  M_{\infty}$.
\end{lemma}
\begin{remark}
    The radius is the distance from $0$ to $z$, which is the same as $\|z\|$.
\end{remark}
We get the radius by taking the intersection of the real axis with a circle in zero with radius $\dist (0,z)$. \ref{Fig.radius}
\begin{proof}
    \proves{lem:construction_radius}
    \leanok
    \uses{lem:Construction.zero_real_root_M_inf, lem:Construction.mul_M_inf, lem:Construction.add_M_Inf, lem:Construction.real_in_M_inf, lem:Construction.im_in_M_inf}  
    We use the fact that $r = \|z\| = \sqrt{\Re(z)^2 + \Im(z)^2}$. Since we have spilt $r$ in already constructed parts, we get that $r \in M_{\infty}$.
\end{proof}

\begin{figure}[h]
    \centering
    \begin{tikzpicture}
        \draw[->] (-0.5,0) -- (4,0) node[right] {$\Re$};
        \draw[->] (0,-0.5) -- (0,4) node[above] {$\Im$};
        \coordinate[label=-135:$0$] (zero) at (0,0);
        \coordinate[label=-135:$1$] (one) at (1,0);
        \coordinate[label=90:$z$] (z) at (3,2);
        \coordinate[label=-135:$r$] (r) at (3.6056,0);
        \fill[black] (zero) circle (2pt);
        \fill[black] (one) circle (1pt);
        \fill[black] (z) circle (2pt);
        \fill[black] (r) circle (2pt);
        \draw (zero) -- (z);
        \draw (zero) [carc=-10:100:3.6056];
    \end{tikzpicture}
    \caption{Construction of $r = \|z\|$}
    \label{Fig.radius}
\end{figure}

\begin{lemma}[Root of pos real number]
    \label{lem:pos_real_root}
    \leanok
    \lean{Construction.one_real_root, Construction.zero_real_root_M_inf}
    If $r \in M_{\infty}\cap\R_{\ge 0}$, then $\sqrt{r}$ is in  $M_{\infty}$.
\end{lemma}

This construction is taken from  \cite{JAN_SCHRÖER:2023}.\\
Draw a circle trough $0$ and $r$ and a line through $1$ parallel to the imaginary axis. Project the intersection to the real axis, using a circle with center $0$ and you get $\sqrt{r}$. (Fig. \ref{Fig.root})
\begin{proof}
    \leanok
    \uses{def:line, def:circle, lem:icc_M_inf, lem:M_in_M_inf} %TODO add rest of dependencies   
Without loss of generality we can assume that $r \ge 1$. 
Otherwise we can use the fact that $\sqrt{r} = \frac{1}{\sqrt{r^{-1}}}$.
The initial step is to define the following lines and circles:
\begin{itemize}
    \item [$l_1$] $:= \{z_1 :=1; z_2 := \imath + 1\}$
    \item [$l_{\Re}$] $:= \{z_1:=0; z_2:= 1\}$
    \item [$c_1$] $:= \{\text{center}:= \frac{r}{2}; \text{radius} := \dist (0,\frac{r}{2})\}$
    \item [$c_2$] $:= \{\text{center}:= 0; \text{radius} := \dist (0,\sqrt{r})\}$
\end{itemize}
It can be observed that both $l_1$ and $l_{\Re}$ are elements of $L(M)$, and that $c_1$ is an element of $C(M)$. 
Furthermore, it can be demonstrated that the $\sqrt{r}$  is an element of $l_{\Re} \cap c_2$. 
Consequently, the only remaining step is to show that $c_2$ is an element of $C(M)$, which is equivalent to proving that there exists a $z$ in $M_{\infty}$ that is also an element of $c_2$. 
This is possible since $0$ is an element of $M_{\infty}$.\\
\underline{Claim}: There exists a $z \in l_1 \cap c_1$, such that $z \in c_2$.
\begin{proof}
    In accordance with the theorem of Pythagoras, it can be demonstrated that$z= 1 \pm \imath \sqrt{r-1}$.
    A further application of the Pythagorean theorem yields the following result: 
    $$\dist (0,z) \overset{z= 1 + \imath \sqrt{r-1}}{=} \sqrt{1^2 + (\sqrt{r-1})^2} = \sqrt{1 + r - 1} = \sqrt{r} = \dist (0,\sqrt{r}).$$
\end{proof}
Therefore, it can be concluded that $\sqrt{r}$ is also constructible.

\end{proof}
\begin{figure}[h]
    \centering
    \begin{tikzpicture}
        \draw[->] (-1,0) -- (7,0) node[right] {$\Re$};
        \draw[->] (0,-1) -- (0,4) node[above] {$\Im$};
        \coordinate[label=-135:$0$] (zero) at (0,0);
        \coordinate[label=-90:$1$] (one) at (1,0);
        \coordinate[label=-45:$r$] (r) at (6,0);
        \coordinate[label=-90:$\frac{r}{2}$] (r2) at (3,0);
        \coordinate[label=135:] (sqrt1) at (1,2.2360679775);
        \coordinate[label=135:$\sqrt{r}$] (sqrt) at (2.4494897428,0);
        \fill[black] (zero) circle (1pt);
        \fill[black] (one) circle (2pt);
        \fill[black] (r) circle (2pt);
        \fill[black] (r2) circle (2pt);
        \fill[black] (sqrt1) circle (2pt);
        \draw (r2) [carc=-10:190:3];
        \draw (zero) [carc=-10:80:2.4494897428];
        \draw (one) -- (sqrt1);
        \fill[red] (sqrt) circle (3pt);
    \end{tikzpicture}
    \caption{Construction of $\sqrt{r}$}
    \label{Fig.root}
\end{figure}

\begin{corollary}[Square root of a complexnumber]
    \label{cor:root_M_inf}
    \uses{def:set_of_constructable_points}
    \leanok
    If $z \in M_{\infty}$, then $\sqrt{z}$ is in  $M_{\infty}$.
\end{corollary}
\begin{proof}
    \leanok
    \uses{lem:construction_radius, lem:pos_real_root, lem:angle_M_inf, lem:Construction.angle_half_M_inf}
    $z = r \cdot e^{\imath \alpha}$ with $r \in \R_{\ge 0}$ and $\alpha \in \R$. Then $\sqrt{z} = \sqrt{r} \cdot e^{\imath \frac{\alpha}{2}}$.
    Now we can use Lemma \ref{lem:construction_radius} and Lemma \ref{lem:pos_real_root} to get that $\sqrt{z} \in M_{\infty}$. For $e^{\imath \frac{\alpha}{2}}$ we can combine Lemma \ref{lem:angle_M_inf} and Lemma \ref{lem:Construction.angle_half_M_inf}. Now we get that $\sqrt{z} \in M_{\infty}$, after Lemma \ref{lem:construction_mul}.

\end{proof}




%2.5
% \begin{lemma}
%   \label{lem:sqrt_in_M_inf}
%   For $M\subseteq \mathbb{C}$ with $0,1 \in M$. If $z$ the $M_{\infty}$, then is $\sqrt{z} \in M_{\infty}$.
% \end{lemma}
% \begin{proof}
%    Let $z = r \exp(\imath \alpha) \in M_{\infty}$, whitout loss of generality we can assume that $r \ne 0$. 
%    Lemma \ref{THM:M_inf_properties}(iv) implies that $r, \exp(\imath \alpha) \in M_{\infty}$. By combining construction \ref{lem:Construction.angle_half_M_inf},
%    which states that $\exp(\imath \alpha / 2) \in M_{\infty}$ and construction \ref{construction_sqrt_pos_real},
%    that gives us that $\sqrt{r} \in M_{\infty}$, we get that $$\sqrt{z} = \sqrt{r} \exp(\imath \alpha / 2) \in M_{\infty}.$$
% \end{proof}


%%FILE blueprint/src/chapter/ch02fieldofconstnum.tex
\chapter{Field of constructable Numbers}
This chapter develops the field structure on $M_{\infty}$ and establishes a set of properties that are utilized to establish a criterion for determining the constructability of a point.
\section[Filed of Constructable Numbers]{Field $M_{\infty}$}
In this section, we will utilise the constructed points from sectoion \ref{basic_constructions} in order to demonstrate that $M_{\infty}$ forms a conjugate (\ref{def:conj_closed}) and quartic (\ref{def:quadritc_closed_field}) closed field.

\begin{theorem}
    \label{thm:MField}
    \leanok
    \lean{MField}
    \uses{def:set_of_constructable_points}
    For $M\subseteq \mathbb{C}$ with $0,1 \in M$. $M_{\infty}$ is a subfield of $\mathbb{C}$.
\end{theorem}
\begin{proof}
    \leanok
    \uses{lem:M_in_M_inf, cor:construction_add, cor:construction_mul_complex, lem:construction_neg, cor:inv_M_inf}
    To show that $M_{\infty}$ is a subfield of $\mathbb{C}$ we have to show that $0,1\in M_{\infty}$ and $M_{\infty}$ is closed under addition, multiplication, subtraction and division.
    \begin{itemize}
        \item [$0,1$:] This follows from $0,1 \in M$ and Lemma \ref{lem:M_in_M_inf}.
        \item [$+$:] For $z_1,z_2 \in M_{\infty}$ we can construct $z_1 + z_2 \in M_{\infty}$. \ref{lem:construction_add}
        \item [$*$:] For $z_1,z_2 \in M_{\infty}$ we can construct $z_1 \cdot z_2 \in M_{\infty}$. \ref{cor:construction_mul_complex}
        \item [$-$:] For $z \in M_{\infty}$ we can construct $-z \in M_{\infty}$. \ref{lem:construction_neg}
        \item [$^{-1}$:] For $z \in M_{\infty}$ with $z \ne 0$ we can construct $z^{-1} \in M_{\infty}$. \ref{cor:inv_M_inf}
    \end{itemize}
\end{proof}

\begin{remark}
   To prove that $M_{\infty}$ is a subfield of $\mathbb{C}$ in Lean we have to create a new instance with carrier $M_{\infty}$.
\end{remark}
\begin{remark}
    \label{rem:MField_Field}
    Since $M_{\infty}$ is a subfield of $\mathbb{C}$, $M_{\infty}$ is a field which is automatically proved in Lean, by \lstinline{infer_instance}.
\end{remark}


\begin{lemma}
    \label{lem:M_inf_properties}
    \leanok
    \lean{MField_i, MField_re_im, MField_re_im', MField_polar}
    \uses{thm:MField}
    For $M\subseteq \mathbb{C}$ with $0,1 \in M$ it holds:
    \begin{itemize}
        \item[(i)] $\imath \in M_{\infty}$.
        \item[(ii)] For $z = x + \imath y \in \mathbb{C}$ the following are equivalent:
            \begin{enumerate}
                \item $z \in M_{\infty}$.
                \item $x, y \in M_{\infty}$.
                \item $x, \imath y \in M_{\infty}$.
            \end{enumerate}
        \item[(iii)] 
            For $0 \ne z = r \exp(\imath \alpha) \in \mathbb{C}$ the following are equivalent:
            \begin{enumerate}
                \item $z \in M_{\infty}$.
                \item $r,\exp(\imath \alpha) \in M_{\infty}$.
            \end{enumerate}
    \end{itemize}
\end{lemma}
\begin{proof}
    \leanok
    \uses{cor:construction_imath, cor:z_iff_re_im_M_inf, cor:construction_imath, cor:construction_radius, lem:angle_M_inf}
    This lemma is a direct consequence of section \ref{basic_constructions}.
    \begin{itemize}
        \item[(i):] We can apply construction \ref{cor:construction_imath}
        \item[(ii):] We can apply construction \ref{cor:z_iff_re_im_M_inf} and \ref{cor:construction_imath}.
        \item[(iii):]  We can apply construction \ref{lem:construction_radius} and \ref{lem:angle_M_inf}.
    \end{itemize}
\end{proof}

\subsection*{quadratic closed}
\begin{definition}[quadratic closed field]
    \label{def:quadritc_closed_field}
    \leanok
    \lean{QuadraticClosed}
    A field $F$ is called \emph{quadratic closed} if for all $x \in F$ there is a $y \in F$ such that $y^2 = x$.
\end{definition}
\begin{remark}
    An equivalent definition is that $F$ is quadratic closed if $F=\{a^2\mid a \in F\}$.
\end{remark}

\begin{lemma}
    \label{lem:M_inf_quad_closed}
    \leanok
    \lean{MField_QuadraticClosed, MField_QuadraticClosed_def}
    \uses{def:set_of_constructable_points, def:quadritc_closed_field}
    For $M\subseteq \mathbb{C}$ with $0,1 \in M$, $M_{\infty}$ is quadratic closed.
\end{lemma}
\begin{proof}
    \uses{thm:MField, cor:root_M_inf}
    \leanok
    It is established that $M_{\infty}$ is a field (see remark \ref{rem:MField_Field}). Furthermore, the corollary \ref{cor:root_M_inf} provides a root $z^{\frac{1}{2}}$ of $z \in M_{\infty}$.    
    $$ z^{\frac{1}{2}} * z^{\frac{1}{2}} = z^{\frac{1}{2}^2} = z^{2\cdot \frac{1}{2}} = z.$$
    Therefore $M_{\infty}$ is quadratic closed.
\end{proof}

\subsection*{Conjugate closed}
\begin{definition}
    \label{def:conj_set}
    \leanok
    \lean{conj_set}
    For a Set $M \subset \mathbb{C}$ we define the \emph{conjugate set} of $M$ as 
    \begin{equation*}
        Conj(M) = \{\overline{z}\mid z\in M\}
    \end{equation*}
\end{definition}

\begin{lemma}
    \label{lem:conj_union}
    \leanok
    \lean{conj_union}
    \uses{def:conj_set}
    For two sets $M,N \subset \mathbb{C}$ $$Conj(M\cup N) = Conj(M) \cup Conj(N).$$
\end{lemma}
\begin{proof}
    \leanok
    For $z \in Conj(M\cup N)$ there is a $w \in M\cup N$ such that $\overline{w} = z$, therefore $ z = \overline{w} \in Conj(M) \cup Conj(N)$. The other direction is analog.
\end{proof}

\begin{lemma}
    \label{lem:conj_conj_id}
    \leanok
    \lean{conj_conj_id}
    \uses{def:conj_set}
    For a set $M \subset \mathbb{C}$ it holds that $Conj(Conj(M)) = M$.
\end{lemma}
\begin{proof}
    \leanok
    $$Conj(Conj(M)) = Conj(\{\overline{z}\mid z\in M\}) = \{\overline{\overline{z}}\mid z\in M\} = \{ z \mid z\in M\} = M.$$
\end{proof}

\begin{definition}
    \label{def:conj_closed}
    \leanok
    \lean{ConjClosed}
    \uses{def:conj_set}
    We call a subset of $\C$ \emph{conjugate closed} if $M= Conj(M)$.
\end{definition}


\begin{lemma}
    \label{lem:conj_MField}
    \leanok
    \lean{MField_conj}
    \uses{def:conj_closed, thm:MField}

    $M_{\infty}$ is conjugate closed.
\end{lemma}
\begin{proof}
    \leanok
    \uses{def:conj_set, lem:construction_conj}
    We can apply construction \ref{lem:construction_conj} and the fact that $\overline{\overline z} = z$ for all $z \in\mathbb{C}$.
\end{proof}

\begin{lemma}
    \label{lem:M_con_M}
    \leanok
    \lean{ConjClosed.M_con_M}
    \uses{def:conj_closed}
    For $M\subseteq \mathbb{C}$ $M\cup Conj(M)$ is conjugate closed.
\end{lemma}
\begin{proof}
    \leanok
    \uses{def:conj_set}
    We can apply Lemma \ref{lem:conj_union} and \ref{lem:conj_conj_id}.
    $$Conj(M\cup Conj(M)) = Conj(M) \cup Conj(Conj(M)) = M \cup Conj(M).$$
\end{proof}

\begin{lemma}
    \label{lem:ConjClosed.Rat_ConjClosed}
    \leanok
    \lean{ConjClosed.Rat_ConjClosed, ConjClosed.Rat_ConjClosed'}
    \uses{def:conj_closed}
    The set of rational numbers is conjugate closed.
\end{lemma}

\begin{proof}
    \leanok
    \uses{def:conj_set}
    For every $r \in \mathbb{Q}$ we have $\overline{r} = r$.
\end{proof}

\begin{lemma}
    \label{lem:ConjClosed.conj_inclusion}
    \leanok
    \lean{ConjClosed.conj_inclusion}
    \uses{def:conj_closed}
    For $M,N \subseteq \mathbb{C}$ with $M \subseteq N$ it holds that $Conj(M) \subseteq Conj(N)$.
\end{lemma}
\begin{proof}
    \leanok
    \uses{def:conj_set}
    For $z \in Conj(M)$ there is a $w \in M$ such that $\overline{w} = z$ and since $M \subseteq N$ we have $w \in N$ and therefore $z \in Conj(N)$.
\end{proof}

\begin{lemma}
    \label{lem:ConjClosed.conj_field}
    \leanok
    \lean{ConjClosed.conj_field}
    \uses{def:conj_closed}
    For a subfield $F$  of $\mathbb{C}$ the conjugate set $Conj(F)$ is a subfield of $\mathbb{C}$.
\end{lemma}
\begin{proof}
    \leanok
    \uses{def:conj_set}
    We have to show that $0,1 \in Conj(F)$ and $Conj(F)$ is closed under addition, multiplication, negation and inversion.
    \begin{itemize}
        \item [$0,1$:] Since $0,1 \in F$ and $\overline{0} = 0, \overline{1} = 1$ we have $0,1 \in Conj(F)$.
        \item [$+$:] For $z_1,z_2 \in F$ we have $\overline{z_1 + z_2} =  \Re(z_1 + z_2)  - \imath\cdot \Im(z_1 + z_2) = \Re(z_1) + \Re(z_2) - \imath\cdot (\Im(z_1) + \Im(z_2)) = \overline{z_1} + \overline{z_2}$. Therefore $Conj(F)$ is closed under addition.
        \item [$-$:] For $z \in F$ we have $\overline{-z} = \Re(-z) - \imath\cdot \Im(-z) = -(\Re(z) - \imath\cdot (\Im(z))) = -\overline{z}$. Therefore $Conj(F)$ is closed under negation.
        \item [$*$:] Since $\overline{z_1 \cdot z_2} = \overline{z_1} \cdot \overline{z_2}$ we have $Conj(F)$ is closed under multiplication.
        \item [$^{-1}$:] Since $\overline{z^{-1}} = \overline{z}^{-1}$ we have $Conj(F)$ is closed under inversion.
    \end{itemize}
\end{proof}

\begin{lemma}
    \label{lem:ConjClosed.re_im_in_L}
    \leanok
    \lean{ConjClosed.ir_L, ConjClosed.im_L}
    \uses{def:conj_closed}
    Let $L$ be a subfield of $\mathbb{C}$, with $L = conj(L)$. For all $z = x + \imath y \in L$ we have $x, \imath y \in L$.
\end{lemma}
\begin{proof}
    \leanok
    Let $z = x + \imath y \in L$. Since $L$ is conjugate closed we know that $\overline{z}=x-\imath y \in L$. This implies
    \begin{equation*}
        \frac{z + \overline{z}}{2} = x \in L
    \end{equation*}
    and therefore also $\imath y = z - x \in L$.
\end{proof}


\begin{lemma}
    \label{lem:ConjClosed.distSq_L}
    \leanok
    \lean{ConjClosed.distSq_L}
    \uses{def:conj_closed}
    Let $L$ be a subfield of $\mathbb{C}$, with $L = conj(L)$, and $z_1, z_2 \in L$.
    For $r := \|z_1-z_2\|$ we get that $r^2 \in L$.
\end{lemma}
\begin{proof}
    \leanok
    \uses{lem:ConjClosed.re_im_in_L}
    For $z_1 = x_1 + \imath y_1$ and $z_2 = x_2 + \imath y_2$ we have
    \begin{equation*}
        r = \|z_1 - z_2\| = \sqrt{(x_1 - x_2)^2 + (y_1 - y_2)^2}
    \end{equation*}
    and therefore
    \begin{equation*}
        r^2 = (x_1 - x_2)^2 + (y_1 - y_2)^2 
    \end{equation*}
    After applying Lemma \ref{lem:ConjClosed.re_im_in_L} we get $r^2 \in L$.
\end{proof}

\begin{lemma}
    \label{lem:ConjClosed.Intersection_line_line}
    \leanok
    \lean{ConjClosed.ill_L, ConjClosed.ill_L'}
    \uses{def:conj_closed}
    Let $L$ be a subfield of $\mathbb{C}$, with $L = conj(L)$. For $i = 1,2,3,4$ let $z_i = x_i + \imath y_i \in L$ with $z_1 \ne z_2$ and $z_3 \ne z_4$. Define
    %TODO: write in two lines
    \begin{equation*}\begin{aligned}
        G_1 := \{\lambda z_1 + (1-\lambda)z_2 \mid \lambda \in \mathbb{R}\},\\
        G_2 := \{\mu z_3 + (1-\mu)z_4 \mid \mu \in \mathbb{R}\}.
    \end{aligned} \end{equation*}
    If $G_1 \cap G_2 \ne \emptyset$ and $G_1 \ne G_2$, it is equivalent 
    \begin{itemize}
        \item $z\in G_1 \cap G_2$.
        \item There are $\lambda, \mu \in \mathbb{R}$ such that:
        \subitem 1. $\lambda(x_1 - x_2)+\mu(x_4-x_3) = x_4-x_2$
        \subitem 2. $\lambda(\imath y_1 - \imath y_2)+\mu(\imath y_4-\imath y_3) = \imath y_4-\imath y_2$
        \subitem 3. $z = \lambda z_1 + (1-\lambda)z_2$
    \end{itemize}
    In this case $z \in L$.
\end{lemma}
\begin{proof}
    \proves{lem:ConjClosed.Intersection_line_line}
    \leanok
    \uses{lem:ConjClosed.re_im_in_L}
    The proof is divided into two parts. Initially, it is demonstrated that $z$ belongs to the intersection of $G_1$ and $G_2$, if and only if there exist real numbers $\lambda, \mu \in \mathbb{R}$, such that the equations $1$, $2$ and $3$ are satisfied.
    Subsequently, it follows that $z$ is an element of $L$.\\
    Equations $1$ and $2$ are equivalent to the following:
    $$ \lambda x_1 + (1 - \lambda)x_2 = \mu x_3 + (1 - \mu)x_4 $$
    $$ \lambda y_1 + (1 - \lambda)y_2 = \mu y_3 + (1 - \mu)y_4 $$ 
    This is the definition of $z \in G_1 \cap G_2$, expressed in terms of its real and imaginary parts.\\
    The third equation is equivalent to $z = \lambda z_1 + (1 - \lambda)z_2$. This allows us to conclude that $z$ belongs to $G_1$ at the point where $G_1$ and $G_2$ intersect. 
    Consequently, we can assume that $z$ belongs to the intersection of $G_1$ and $G_2$.

    Now we can show that $z$ is an element of $L$.\\
    Since we know that z is equal to $\lambda z_1 + (1-\lambda)z_2$ and $z_1, z_2 \in L$ we only have to show that $\lambda \in L$. Here for we use the equations from the first part of the proof.
    \begin{align*}
        \RNum{1} &&\lambda(x_1 - x_2)+\mu(x_4-x_3) &= x_4-x_2\\
        \RNum{2} &&\lambda(\imath y_1 - \imath y_2)+\mu(\imath y_4-\imath y_3) &= \imath y_4-\imath y_2
    \end{align*}
    Now we solve \RNum{2} for $\mu$ 
    \begin{align*}
        && \lambda(\imath y_1 - \imath y_2)+\mu(\imath y_4-\imath y_3) &= \imath y_4-\imath y_2 && \mid -\lambda(\imath y_1 - \imath y_2)\\
        &\Leftrightarrow & \mu(\imath y_4-\imath y_3) &= \imath y_4-\imath y_2 - \lambda(\imath y_1 - \imath y_2) && \mid \div \imath(y_4-y_3)\\
        &\Leftrightarrow & \mu &= \frac{\imath y_4-\imath y_2 - \lambda(\imath y_1 - \imath y_2)}{\imath y_4-\imath y_3}\\
       %&\Leftrightarrow & \mu(y_4-y_3) &= y_4-y_2 - \lambda(y_1 - y_2) &&\mid \div \imath(y_4-y_3)\\
        %&\Leftrightarrow & \mu &= \frac{y_4-y_2 - \lambda(y_1 - y_2)}{y_4-y_3}\\
    \end{align*}
    Since we divided by $\imath (y_4-y_3)$ we need to assume that $y_4 \ne y_3$, so we need to first handle the case $y_4 = y_3$.\\
    If $y_4 = y_3$ we have $\lambda(\imath y_1 - \imath y_2) = \imath y_4-\imath y_2$ and since $y_4 = y_3$ $y_1 \ne y_2$, because otherwise both Lines would be parallel to the real line and therefore $G_1 = G_2$ or $G_1 \cap G_2 = \varnothing$. 
    %Todo: Remove + 0 after the issue is solved
    Therefore $\lambda = \frac{\imath y_4-\imath y_2}{\imath y_1 - \imath y_2}$. Using the fact that real part and the imaginary part times $\imath + 0$ are in $L$ \ref{lem:ConjClosed.re_im_in_L} we have written $\lambda$ as a fraction of two elements in $L$. 
    It can thus be concluded that $\lambda$ is in $L$, which implies that $z = \lambda z_1 + (1-\lambda)z_2$ is in $L$.\\
   
    Now we insert $\mu$ in \RNum{1} and solve for $\lambda$.\\
    %TODO: Change block such that it will be also shown in the Web version 
    \resizebox{1\linewidth}{!}{
    \begin{minipage}{\linewidth}
    \begin{alignat*}{3}
        && \lambda(x_1 - x_2)+\mu(x_4-x_3) &= x_4-x_2 && \mid \RNum{1} \leftarrow \RNum{2}\\
        &\Leftrightarrow & \lambda(x_1 - x_2)+\frac{\imath y_4-\imath y_2 - \lambda(\imath y_1 - \imath y_2)}{\imath y_4-\imath y_3}(x_4-x_3) &= x_4-x_2 && \mid \cdot (\imath y_4-\imath y_3)\\
        &\Leftrightarrow & \lambda(x_1 - x_2)(\imath y_4-\imath y_3)+(\imath y_4-\imath y_2 - \lambda(\imath y_1 - \imath y_2))(x_4-x_3) &= (x_4-x_2)(\imath y_4-\imath y_3) && \mid - (x_1 - x_2)(\imath y_4-\imath y_3)\\
        %&\Leftrightarrow & \lambda(y_4-y_3)(x_1 - x_2)+(y_4-y_2)(x_4-x_3) - \lambda(y_1 - y_2)(x_4-x_3) &= (x_4-x_2)(y_4-y_3) && 
        &\Leftrightarrow & \lambda((x_1 - x_2)(\imath y_4-\imath y_3) - (\imath y_4-\imath y_2)(x_4-x_3)) &= (x_4-x_2)(\imath y_4-\imath y_3) - (\imath y_4-\imath y_2)(x_4-x_3) && \mid \div ((\imath y_4-\imath y_3)(x_1 - x_2)- (\imath y_1 - \imath y_2)(x_4-x_3))\\
        %&\Leftrightarrow & \lambda((y_4-y_3)(x_1 - x_2)- (y_1 - y_2)(x_4-x_3)) &= (x_4-x_2)(y_4-y_3) - (y_4-y_2)(x_4-x_3) && \mid \div ((y_4-y_3)(x_1 - x_2)- (y_1 - y_2)(x_4-x_3))\\
        &\Leftrightarrow & \lambda &= \frac{(x_4-x_2)(\imath y_4-\imath y_3) - (\imath y_4-\imath y_2)(x_4-x_3)}{(\imath y_4-\imath y_3)(x_1 - x_2)- (\imath y_1 - \imath y_2)(x_4-x_3)}\\
    \end{alignat*}
    \end{minipage}
    }\\
    We need to check that the denominator $(y_4-y_3)(x_1 - x_2)- (y_1 - y_2)(x_4-x_3)$ is not zero.
    Assume that its = then we would have $(y_4-y_3)(x_1 - x_2) = (y_1 - y_2)(x_4-x_3)$, which is equivalent to $\frac{y_4-y_3}{x_4-x_3} = \frac{y_1 - y_2}{x_1 - x_2}$. 
    This would mean that the two lines are parallel and therefore $G_1 = G_2$ or $G_1 \cap G_2 = \varnothing$.\\
    Thus we can assume that the denominator is not zero and therefore we can write $\lambda$ as a fraction of two elements in $L$. Therefore $\lambda$ is in $L$, wich implies that $z = \lambda z_1 + (1-\lambda)z_2$ is in $L$.
\end{proof}

%3.17
\begin{lemma}
    \label{lem:ConjClosed.Intersection_line_circle}
    \leanok
    \uses{def:conj_closed}
    Let $L$ be a subfield of $\mathbb{C}$, with $L = conj(L)$. For $i = 1,2,3$ let $z_i = x_i + \imath y_i \in L$ with $z_1 \ne z_2$, and let $r > 0$ in $\mathbb{R}$ with $r^2 \in L$. Define
    \begin{equation*}\begin{aligned}
        G := \{\lambda z_1 + (1-\lambda)z_2 \mid \lambda \in \mathbb{R}\},\\
        C := \{z = x + \imath y \in \mathbb{C} \mid \|z - z_3\| = r\}.
    \end{aligned} \end{equation*}
    Assume $G \cap C \ne \emptyset$; then the following are equivalent:
    \begin{itemize}
        \item $z\in G \cap C$.
        \item There is a $\lambda \in \mathbb{R}$ with $\lambda^2 a+ \lambda b + c = 0$ where
        \begin{align*}
            a &:= (x_1 - x_2)^2 + (\imath y_1 - \imath y_2)^2,\\
            b &:= 2(x_1 - x_2)(x_2 - x_3) - 2(\imath y_1 - \imath y_2)(\imath y_2 - \imath y_3),\\
            c &:= (x_2 - x_3)^2 + (\imath y_2 - \imath y_3)^2 - r^2,
        \end{align*}
        and $z = \lambda z_1 + (1-\lambda)z_2$.
    \end{itemize}
    In this case $z \in L(\sqrt{w})$ for an $w \in L$.
\end{lemma}

\begin{proof}
    \leanok
    \uses{lem:ConjClosed.distSq_L, lem:ConjClosed.re_im_in_L}
First we have to show $z \in G \cap C$ iff and only iff there exists a $\lambda \in \mathbb{R}$ with $\lambda^2 a+ \lambda b + c = 0$ and $z = \lambda z_1 + (1-\lambda)z_2$.

"$\Rightarrow:$"  If z belongs to the intersection of G and C, then z satisfies the equations of C and G. Consequently
\begin{align*}
    &&\|z - z_3\| &= r &&\mid 0\le \|\cdot\|\text{ and }0\le r\\
    &\Leftrightarrow& \|z - z_3\|^2 &= r^2\\
    &\Leftrightarrow& (x - x_3)^2 + (\imath y - \imath y_3)^2 &= r^2 &&\mid x = \lambda x_1 - \lambda x_2 + x_2 \text{ and }\\ 
    &&&&&\mid y = \lambda y_1 - \lambda y_2 + y_2\\
    &\Leftrightarrow& (\lambda x_1 - \lambda x_2 + x_2 - x_3)^2 + \\
    &&(\imath (\lambda y_1 - \lambda y_2 + y_2 - y_3))^2 &= r^2\\
    &\Leftrightarrow& \lambda ^2 ((x_1 - x_2)^2 + (\imath y_1 - \imath y_2)^2) + \\
    &&\lambda (2(x_1 - x_2)(x_2 - x_3) - 2(\imath y_1 - \imath y_2)(\imath y_2 - \imath y_3)) + \\
    &&(x_2 - x_3)^2 + (\imath y_2 - \imath y_3)^2 &= r^2\\
\end{align*}
"$\Leftarrow:$" Since $z = \lambda z_1 + (1-\lambda)z_2$ we get $z$ in $G$ and can use the equations from the first part of the proof to show that $z$ is in $C$.

Now we can show that there exists a $w \in L$ such that $z \in L(\sqrt{w})$.\\
Since we know that $z = \lambda z_1 + (1-\lambda)z_2$ and $z_1, z_2 \in L$ we only have to show that $\lambda \in L(\sqrt{w})$.
To do this, we use the equations from the first part of the proof. 
Since $\lambda$ is a solution of a quadratic equation we now get that $\lambda$ is equal to $\frac{-b \pm \sqrt{b^2 - 4ac}}{2a}$. Since $a,b,c \in L$ we get $w = b^2 - 4ac \in L$ so $\lambda \in L(\sqrt{w})$. Therefore $z = \lambda z_1 + (1-\lambda)z_2$ is in $L(\sqrt{w})$.
%TODO: add proof
\end{proof}
%3.18
\begin{lemma} 
    \label{lem:ConjClosed.Intersection_circle_circle}
    \lean{ConjClosed.icc_L, ConjClosed.icc_L'}
    \leanok
    \uses{def:conj_closed}
    Let $L$ be a subfield of $\mathbb{C}$, with $L = conj(L)$. For $i = 1,2 $ let $z_i = x_i + \imath y_i \in L$ with $z_1 \ne z_2$ and let $r_i > 0$ in $\mathbb{R}$ with $r_i^2 \in L$. Define
    \begin{equation*} \begin{aligned}
        C_1 := \{z = x + \imath y \in \mathbb{C} \mid \|z - z_1\| = r_1\},\\
        C_2 := \{z = x + \imath y \in \mathbb{C} \mid \|z - z_2\| = r_2\}.
    \end{aligned} \end{equation*}
    Assume $C_1 \cap C_2 \ne \emptyset$ and $C_1 \ne C_2$. Then  ther  exists $z_1, z_2 \in L$ such that
     $$G := \{ \lambda z_1 + (1-\lambda) z_2 \mid \lambda \in \R \}$$ % \{x+\imath y \in \mathbb{C} \mid 2(x_2 - x_1)x - 2(\imath y_2 - \imath y_1)\imath y = r_1^2 - r_2^2 + x_2^2 - x_1^2 + (\imath y_2)^2 - (\imath y_1)^2\} $$
    is a real line, and $$ C_1 \cap C_2 = G \cap C_1 = G \cap C_2. $$
    For $z \in C_1 \cap C_2$ there is a $w \in L$ such that $z \in L(\sqrt{w})$.
\end{lemma}
\begin{proof}
    \uses{lem:ConjClosed.distSq, lem:ConjClosed.Intersection_line_circle}
    The initial step is to demonstrate that $ z\in C_1 \cap C_2 \Leftrightarrow \exists G,  z\in G \cap C_1 \land z \in G \cap C_2$.\\
    "$\Rightarrow$": If $z$ is to be in both $G \cap C_1$ and $G \cap C_2$, it must be the case that $z$ is in both $C_1$ and $C_2$. Consequently, it is also in $C_1 \cap C_2$. \\
    "$\Leftarrow$": We begin by establishing that $z=x + iy\in C_1\cap C_2$ is equivalent to:
        $$\|z-z_1\| = r_1 \land \|z - z_2\| = r_2 \Leftrightarrow 2(x_2-x_1)x-2(\imath y_2-\imath y_1)\imath y = r_1^2-r_2^2-x_1^2+x_2^2+(\imath y_1)^2-(\imath y_2)^2$$
        The remaining task is to identify two elements in $L$ that satisfy the given equation. This can be achieved by considering three cases:
        In the initial case, where $x_1 = x_2$, it can be demonstrated that $y_1 \ne y_2$, as otherwise it would follow that $C_1=C_2$. Here, we can choose
        \begin{align*}
            z_1 &:= 1 + \imath (\frac{r_1^2 - r_2^2 + (\imath y_2)^2 - (\imath y_1)^2}{-2(y_2-y_1)})\\
            z_2 &:= 0 + \imath (\frac{r_1^2 - r_2^2 + (\imath y_2)^2 - (\imath y_1)^2}{-2(y_2-y_1)})
        \end{align*}
        In the second case we have $y_1 = y_2$ and $x_1 \ne x_2$. Here we can choose 
        \footnote[1]{During the process of formalising this proof, it became evident that for $y_1=y_2=0$ and $L\subseteq\R$, there does not exist a line $G$ with the property that $G\cap C_1=G\cap C_2=C_1\cap C_2$.
            This error can also be found in the source \cite{JAN_SCHRÖER:2023} and was noticed too late, so it could not be corrected in time. The existence of an $w$ such that $z\in L(\sqrt{w})$ is still correct and is the result we are interested in.}
        \begin{align*}
            z_1 &:= (\frac{r_1^2 - r_2^2 + x_2^2 - x_1^2}{2(x_2-x_1)}) + \imath y_1\\
            z_2 &:= (\frac{r_1^2 - r_2^2 + x_2^2 - x_1^2}{2(x_2-x_1)}) - \imath y_1
        \end{align*}
    %$$\|z - z_1\| = r_1 \land \|z - z_2\| = r_2 \Leftrightarrow 2(x_2-x_1)x-2(\imath y_2-\imath y_1)\imath y = r_1^2-r_2^2-x_^2+x_2^2+(\imath y_1)^2-(\imath y_2)^2$$ % \overset{0\le\|\cdot\|\land0\le r}{\Leftrightarrow}
        For $ x_1 \ne x_2$ and $y_1 \ne y_2$ chose 
        \begin{align*}
            z_1 &:= (\frac{r_1^2 - r_2^2 + x_2^2 - x_1^2 + (\imath y_2)^2 - (\imath y_1)^2+2(\imath (y_2)-(\imath y_1))(\imath y_1)}{2(x_2-x_1)}) + \imath y_1\\ % x_1 + \imath (\frac{r_1^2 - r_2^2 + x_2^2 - x_1^2 + (\imath y_2)^2 - (\imath y_1)^2-2(x_2-x_1)x_1}{2(y_2-y_1)})\\
            z_2 &:= (\frac{r_1^2 - r_2^2 + x_2^2 - x_1^2 + (\imath y_2)^2 - (\imath y_1)^2+2(\imath (y_2)-(\imath y_1))(\imath y_2)}{2(x_2-x_1)}) + \imath y_2\\
        \end{align*}
    
        Since the points $z_1$ and $z_2$ lie in $L$, we can conclude that the line $G$ lies in the set of lines of $L$. This allows us to apply the results stated in lemma \ref{lem:ConjClosed.Intersection_line_circle} to obtain $w$, with the result that $z$ is contained within the set $L(w)$.
\end{proof}

\section[K zero]{The Field $\mathcal{K}_0$}
This section develops a conjugation-closed field that depends on the set $\mathcal{M}$.
\begin{definition}
    \label{def:K_M_0}
    \lean{K_zero}
    \leanok
    Let $\mathcal(M)\subseteq\mathbb{C}$ with $0,1 \in \mathcal{M}$. Define:
    \begin{equation*}
        K_0 := \mathbb{Q}(\mathcal{M}\cup Conj(\mathcal{M}))
    \end{equation*}
\end{definition}

\begin{lemma}
    \label{lem:conj_adjion}
    \leanok
    \lean{conj_adjion, conj_adjion'}
    \uses{def:K_M_0}
    Let $K$ be an conjugation closed intermediate field of $\Q$ and $\C$ and $ M \subset \mathbb{C}$ be a subset with $M = conj(M)$. Then
    $K(M)$ is conjugate closed.
\end{lemma}
\begin{proof}
    \leanok
    \uses{lem:ConjClosed.conj_inclusion, lem:ConjClosed.conj_field}
    In reference \ref{lem:ConjClosed.conj_field}, it was demonstrated that for a field F, the field of complex numbers, $Conj(F)$ is a field. 
    It can thus be concluded that $Conj(K(M))$ is also a field.
    As both $K$ and $M$ are subsets of $K(M)$, it can be inferred from lemma \ref{lem:ConjClosed.conj_inclusion} that $Conj(K) \overset{ConjClosed}{=} K$ and $Conj(M) \overset{ConjClosed}{=} M$ are subsets of $Conj(K(M))$. 
    As $K(M)$ is the smallest subfield of $\C$ that includes $K$ and $M$, it can be concluded that $$K(M) \subseteq Conj(K(M)).$$
    Furthermore, if we apply $Conj$ to both sides and again infer \ref{lem:ConjClosed.conj_inclusion}, we obtain $$Conj(K(M)) \subseteq Conj(Conj(K(M))) = K(M),$$ 
    which leads to the conclusion that $Conj(K(M)) = K(M)$.
\end{proof}

\begin{corollary}
    \label{cor:K_zero_conjectClosed}
    \leanok
    \lean{K_zero_conjectClosed}
    \uses{def:K_M_0}
    For $M\subseteq \mathbb{C}$ with $0,1 \in M$, $K_0$ is conjugate closed.
\end{corollary}
\begin{proof}
    \leanok
    \uses{lem:conj_adjion, lem:ConjClosed.Rat_ConjClosed, lem:M_con_M}
    By employing the preceding lemma, it is sufficient to demonstrate that $\Q$ and $M \cup Conj(M)$ are conjugate closed,
    which can be inferred from \ref{lem:ConjClosed.Rat_ConjClosed} and \ref{lem:M_con_M}
\end{proof}

\begin{lemma}
    \label{lem:K_zero_in_MField}
    \leanok
    \lean{K_zero_in_MField}
    For $M\subseteq \mathbb{C}$ with $0,1 \in M$ it holds that $K_0 \subseteq M_{\infty}$.
\end{lemma}
\begin{proof}
    \leanok
    \uses{lem:M_M_inf, lem:conj_M_Inf, def:K_M_0}
From the definition of $K_0 := \Q(M\cup Conj(M))$, 
it can be seen that this is the smallest subfield of $\C$ containing both $Q$ and $M\cup Conj(M)$. 
Consequently, it is sufficient to demonstrate that both $\Q$ and $M\cup Conj(M)$ are contained within $M_{\infty}$. 
Since $\Q$ is contained in every subfield of $\C$, it is therefore also contained in $M_{\infty}$. 
Furthermore, since $M$ is contained in $M_{\infty}$ (see \ref{lem:M_in_M_inf}) and $M$ is conjugate closed (\ref{lem:conj_MField}), we can conclude that $M\cup Conj(M) \subseteq \Q$.
\end{proof}

\section{Classification of Constructable Numbers}
The following section will demonstrate that for an element to be constructible, the degree over $K_0$ must be equal to $2^m$ for some natural number $m$.
%The following section will employ the following lemma, which has not yet been proven in Lean.

\begin{lemma}
    \label{lem:dergree_two_eq_sqr}
    \lean{dergree_two_eq_sqr}
    \leanok
    %For to abetory fields $E$ and $F$ with
    Let $K, L$ be subfield of $\C$, with $K\le L$. Then $[L:K] = 2$ is equivalent to the existence of a $w*w \in K$ with $w \notin K$ and $L = K(w)$.
\end{lemma}
\begin{proof}
    $(ii)\implies (i)$: Let $w$ be as in $(ii)$.Then $\sqrt{w}$ is a root of $X^2 - w \in K[X]$. Since $\sqrt{w} \notin K$ this polynomial is irreducible in $K[X]$. Therefore $[L:K] = 2$.\\
    $(i)\implies (ii)$: Let $[L:K] = 2$ and $\alpha \in L \setminus K$. Then $K(\alpha) = L$ and 
    $$\mu_{\alpha, K}=X^2 + bX + c \quad b,c \in K$$
    This implies 
    $$\alpha = -\frac{b}{2} \pm \sqrt{\frac{b^2}{4} - c} $$
    Now let $w := \frac{b^2}{4} - c \in K$ then we get $L = K(\alpha) = K(\sqrt{w})$.
\end{proof}


\begin{lemma}
    \label{lem:Z_in_Minf_imp_eq} 
    \leanok
    \lean{chain_iff_adjion_roots}
    For $z \in \mathbb{C}$
    there exists an $n\ge 0$ and a chain 
    $$K_0 = L_1 \subset L_2 \subset \ldots \subset L_n \subset \mathbb{C}$$
    of subfields of $\mathbb{C}$ such that $z \in L_n$ and 
    $ [L_{i+1}:L_i] = 2 $ for $i = 0, \ldots, n-1$. This is equivalent to:\\
    There is an intermediate field $L$ of $\mathbb{C}/K_0$ with $z \in L$, so that $L$ arises from $K_0$ by a finite sequence of adjunctions of square roots.
\end{lemma}
\begin{proof}
    \uses{lem:dergree_two_eq_sqr}
    "$\Rightarrow:$" Let $n$ be the smallest natural number such that there exists a chain of subfields $L_1 \subset L_2 \subset \ldots \subset L_n$ with $z \in L_n$ and $[L_{i+1}:L_i] = 2$ for $i = 0, \ldots, n-1$. Now lemma \ref{lem:dergree_two_eq_sqr} gives us that there exists a $w \in L_{n-1}$ such that $L_n$ arises from $L_{n-1}$ by adjoining $\sqrt{w}$.\\
    "$\Leftarrow:$" Let $L$ be an intermediate field of $\mathbb{C}/K_0$ with $z \in L$ and $L$ arises from $K_0$ by a finte sequence of adjunctions of square roots. Then there exists a chain of subfields $L_1 \subset L_2 \subset \ldots \subset L_n$ with $z \in L_n$ and $L_{i+1}:= L_i(\sqrt{w_i})$ for $i = 0, \ldots, n-1$.
\end{proof}

\begin{lemma}
    \label{lem:Mi_chain}
    \leanok
    \lean{M_i_in_K_i}
    For $M_n$ exists a chain of intermediate fields $K_0 \le K_1 \le \ldots \le K_n$ such that $M_i\subset K_i$ and $K_i:= K_{i-1}(X_i)$ for a set of square roots $X_i$ of elements of $K_{i-1}$.
\end{lemma}

\begin{proof}
\leanok
\uses{cor:K_zero_conjectClosed, lem:ConjClosed.Intersection_line_line, lem:ConjClosed.Intersection_line_circle, lem:ConjClosed.Intersection_circle_circle, lem:ConjClosed.conj_inclusion}
Induction over $n$
\begin{itemize}
    \item Base case $n=1$: \\
        $K_0 \le K_0$ and $K_0$ is conjugation closed \ref{cor:K_zero_conjectClosed}.
    \item induction hypothesis: \\
        Assume that for $n$ there is a chain of conjugation closed  intermediate fields $K_0 \le K_1 \le \ldots \le K_n$ such that $M_i\subset K_i$ and $K_i:= K_{i-1}(X_i)$ for a set of square roots $X_i$ of elements of $K_{i-1}$.
    \item Inductive step $n \to n+1$: \\
        For $z\in M_{n+1}$ there are four cases: 
        \begin{itemize}
            \item [$z \in M_n$:] By the induction hypothesis $z \in K_n$ and $K_n$ is conjugation closed and arises from $K_0$ by a sequence of adjunctinos of square roots.
            \item [$z \in ILL (M_n)$:] By the induction hypothesis $z \in ILL (K_n)$ and using \ref{lem:ConjClosed.Intersection_line_line} we get that $z \in K_{n}$.
            \item [$z \in ILC (M_n)$:] By the induction hypothesis $z \in ILC (K_n)$ and using \ref{lem:ConjClosed.Intersection_line_circle} there is a $w \in K_n$ such that $z \in K_n(\sqrt{w})$ insert  $\sqrt{w}, \overline{\sqrt{w}}$ to $X_n$.
            \item [$z \in ICC (M_n)$:] By the induction hypothesis $z \in ICC (K_n)$ and using \ref{lem:ConjClosed.Intersection_circle_circle} there is a $w \in K_n$ such that $z \in K_n(\sqrt{w})$ insert $\sqrt{w}, \overline{\sqrt{w}}$ to $X_n$.
        \end{itemize}
\end{itemize}
\end{proof}

\begin{theorem}[constructable iff chain dergee2]
    \label{thm:Classfication_z_in_M_inf}
    \lean{Classfication_z_in_M_inf, adjoin_in_MField'}
    \leanok
    For $z \in \mathbb{C}$, $z \in M_{\infty}$ is equivalent to:\\
    There is a $0\le n$ and a chain 
    $$K_0 = L_1 \subset L_2 \subset \ldots \subset L_n \subset \mathbb{C}$$
    of subfields of $\mathbb{C}$ such that $z \in L_n$ and 
    $$ [L_{i+1}:L_i] = 2 \quad \text{for} \quad i = 0, \ldots, n-1.$$
\end{theorem}

\begin{proof}
    \leanok
    \uses{lem:Mi_chain, lem:Z_in_Minf_imp_eq, lem:M_inf_quad_closed}
    "$\Leftarrow:$"
    It can be shown by induction that $L_n$ is contained within $M_{\infty}$. 
    Therefore, it can be inferred that $z$ is also contained within $M_{\infty}$.
    \begin{itemize}
        \item Base case $n=1$: \\
            $L_1 = K_0 \subseteq M_{\infty}$ \ref{lem:K_zero_in_MField}.
        
        \item induction hypothesis: \\
            Assume that for $n$: $\forall i < n: L_i \le L_{i+1} \land [L_{i+1}:L_i]=2$ implies $L_n \subseteq M_{\infty}$.
        \item Inductive step $n \to n+1$: \\
            Given that $[L_{n+1}:L_n] = 2$, it follows from the conclusions of Lemma \ref{lem:dergree_two_eq_sqr} that there exists a $w \in L_n$ with the property that $ \sqrt{w} \notin L_n$ and $L_{n+1} = L_n(w)$.
            By the induction hypothesis, it can be inferred that $L_n \subseteq M_{\infty}$. Since $w \in L_n  \subseteq M_{\infty}$ and $ M_{\infty}$ is quadratic closed (\ref{lem:M_inf_quad_closed}) $L_n(w) = L_{n+1} \subseteq M_{\infty}$.
    \end{itemize}
    "$\Rightarrow:$" There exists a $n$ such that $z\in M_n$, and we know that there exists a $K_n$ with $M_n \subseteq K_n$ which is derived from $K_0$ by successive adjoing square roots \ref{lem:Mi_chain}. 
    We can conclude that there is a $K$, which is derived from $K_0$ by successive adjoining square roots, and that $z\in K$. 
    Since $M_i$ is finite, we get that we adjoin finitely many square roots and so we evoke \ref{lem:Z_in_Minf_imp_eq}. 

\end{proof}

\begin{lemma}
    \label{lem:Classfication_z_in_M_inf_2m}
    \lean{Classfication_z_in_M_inf_2m}
    \leanok
    For $z \in \mathbb{C}$, $z \in M_{\infty}$ implies there exists a $m$ such that $[z:K_0] = 2^m$.
\end{lemma}

\begin{proof}
    \leanok
    \uses{thm:Classfication_z_in_M_inf}
    By Theorem \ref{thm:Classfication_z_in_M_inf}, it can be inferred that there exists a chain of subfields, $K_0 = L_1 \subset L_2 \subset \dots \subset L_n \subset \mathbb{C}$, with $z \in L_n$ and $[L_i : L_{i+1}] = 2$ for $i = 0, 1, \dots, n-1$. 
    Moreover, the multiplicativity formula for degrees indicates that the degree of the extension $[L_n : K_0]$ is equal to the product of the degrees of the extensions $[L_n : L_{n-1}] \cdot [L_{n-1} : L_{n-2}] \cdot \dots [L_2 : L_1]$.
    Thus, we have that $[L_n : K_0] = 2^n$. 
    The fact that $z\in L_n$ implies that $K_0(z)\subseteq L_n$. 
    It thus follows that the index of the field extension $[L_n : K_0] = [L_n : K_0(z)]\cdot [K_0(z) : K_0]$, which implies that $[K_0(z) : K_0]$ is a divisor of $2^n$. 
\end{proof}

\begin{corollary}
    \label{cor:Classfication_z_in_M_inf_2m_not}
    \leanok
    \lean{Classfication_z_in_M_inf_2m_not}
    For $z \in \mathbb{C}$, if there is no $m$ such that $[z:K_0] = 2^m$ then $z \notin M_{\infty}$.
\end{corollary}

\begin{proof}
    \leanok
    \uses{lem:Classfication_z_in_M_inf_2m, lem:dergree_two_eq_sqr}
    Contraposition of Lemma \ref{lem:Classfication_z_in_M_inf_2m}.
\end{proof}

\begin{corollary}
    \label{cor:Classfication_z_in_M_inf_2m_rat}
    \lean{Classfication_z_in_M_inf_2m_rat}
    \leanok
    For $z \in \mathbb{C}$, $z \in M_{\infty}$ with $[K_0:\Q] = 2^n$,  if there is no $m$ such that $[z:\Q] = 2^m$ then $z \notin M_{\infty}$.
\end{corollary}

\begin{proof}
    \leanok
    \uses{cor:Classfication_z_in_M_inf_2m_not}
    A combination of the multiplicativity formula for degrees and corollary \ref{cor:Classfication_z_in_M_inf_2m_not}.
\end{proof}

%%FILE blueprint/src/chapter/ch0xancientconstructionProplems.tex
\chapter{Ancient Construction Problems}
This chapter will employ the results to demonstrate the impossibility of trisecting the angle and doubling the cube. 
This formalisation is based on the work conducted during my project in Bonn during the Lean Course WiSe 23/24: 
\url{https://github.com/Louis-Le-Grand/LeanCourse23Fork/tree/master/LeanCourse/Project}

\section{Doubling the cube}
The doubling of the cube, also known as the Delian problem, represents an ancient geometric problem.
The objective is to construct the edge of a second cube whose volume is double that of the first, using only a ruler and compass, given the edge of a cube.
The construction of a second cube with double the volume of the original cube begins with a cube of volume $a^3$, where $a$ is the length of an edge. 
Thus, a cube with double the volume ($2\cdot a^3$) has an edge length of the cube root of two times the length of the original edge.
If we now take the unit cube and reduce $\mathcal{M}$, the problem is as follows:
\begin{problem}
    Let $\mathcal{M} = \{0,1\}$.  
    $$\text{Is }\sqrt[3]{2} \overset{?}{\in} \mathcal{M}_{\infty}?$$
\end{problem}

\begin{lemma}($\sqrt[3]{2}$ is irrational)
    \label{lem:irrational_thirdroot_two}
    \leanok
    \lean{irrational_thirdroot_two}
    The third root of $2$ is an irrational number.
\end{lemma}
\begin{proof}
    \leanok
    The following theorem will be used without proof, as it is already available in MathLib:
    \begin{theorem*}
        For any $x\in(\R\setminus\Z)$ if there exist an $n\in \N_{>0}$ and $m\in\Z$ such that $m = x^n$, then $x$ is rational. 
    \end{theorem*}
    The fact that $(\sqrt[3]{2})^3=2$, allows us to deduce that the only remaining task is to prove that it  is not an integer. 
    This can be observed trough two relations.
    \begin{align}
        2^{\frac{1}{3}} &< 2 \\
        2^{\frac{1}{3}} &> 1
    \end{align}
\end{proof}

\begin{lemma}
    $P := X^3 - 2$ is irreducible over $\mathbb{Q}$.
\label{lem:irreducible_over_Q}
\leanok
\lean{P_irreducible}
\end{lemma}
\begin{proof}
    \leanok
    \uses{lem:irrational_thirdroot_two}
    Since $\mathbb{Q}$ is $a$ subfield of $\mathbb{C} [X]$, we know that
    \begin{equation*}
        X^3 - 2 = (X - \sqrt[3]{2})(X -\zeta_3 \sqrt[3]{2})(X -\zeta_3^2 \sqrt[3]{2})
    \end{equation*}
    Suppose $P$ is reducible, then
    \begin{equation*}
        X^3 - 2 = (X - a)(X^2 + bX + c)\text{, with } a, b, c \in \mathbb{Q}
    \end{equation*}
    In particular it has a zero in $\mathbb{Q}$, so there is a rational number $a$ such that $a^3 = 2$.\newline
    But we know that $\zeta_3 \sqrt[3]{2}$ and $\zeta_3^2 \sqrt[3]{2}$ are not real numbers and $\sqrt[3]{2}$ is not rational \ref{lem:irrational_thirdroot_two}.
    So $P$ is irreducible over $\mathbb{Q}$.
\end{proof}

\begin{theorem}
    \label{thm:third_root_of_two_not_in_M_inf}
    \leanok
    \lean{third_root_of_two_not_in_M_inf}
    The cube can't be doubled using a compass and straightedge.
\end{theorem}
\begin{proof}
    \leanok
    \uses{lem:irrational_thirdroot_two, lem:irreducible_over_Q, cor:Classfication_z_in_M_inf_2m_not}
    By applying the corollary \ref{cor:Classfication_z_in_M_inf_2m_not}, it is sufficient to proof that no $m\in \N$ exists such that
    $$2^m\overset{?}{=}[\sqrt[3]{2}:\Q({0,1})]\overset{0,1\in\Q}{=} [\sqrt[3]{2}:\Q] = \text{degree}(\mu_{\Q,\sqrt[3]{2}}).$$
    Since $P$ is irreducible over $\mathbb{Q}$ \ref{thm:irreducible_over_Q}, monic and has $\sqrt[3]{2}$ as a zero, we know that $[\mathbb{Q}(\sqrt[3]{2}):\mathbb{Q}] = 3$.
    And since $$ 3  \equiv_2 1 \neq 0  \equiv_2 = 2^m \qquad\forall m \in \N$$
    we can conclude that the cube can't be doubled using a compass and straightedge.
    
\end{proof}

\section{Trisection of an angle}
The trisection of an angle with a compass and ruler can be reduced to the following problem:\\
Let $\mathcal{M} = \{a, b, c\}$ with $a, b, c$ not on a line and $\alpha := \angle (b - a, c - a)$ be the resulting angle.
Then $\alpha$ can be trisected if and only if there is a point $d\in \mathcal{M}_{\infty}$ such that $\angle (b - a, d - a) = \alpha/3$. 
The use of a normed set $\mathcal{M} = \{0,1,e^{i\alpha}\}$ leads to the following problem:
\begin{problem}
    Let $\mathcal{M} = \{0,1,\exp(\textbf{i} \alpha)\}$. $$\text{Is }\exp(\textbf{i} \alpha/3) \overset{?}{\in} \mathcal{M}_{\infty}?$$
\end{problem}
In this context, since the numbers zero and one are rational numbers, it can be concluded that $K_0$ is equal to
$$ K_0 = \mathbb{Q}(\mathcal{M}\cup Conj(\mathcal{M})) = \mathbb{Q}(e^{\imath\alpha},\overline{e^{\imath\alpha}}). $$
Given the corollary reference and the fact that $e^{\imath\alpha/3}$ is a zero of $X^3 - e^{\imath\alpha}\in \mathbb{Q}(e^{\imath\alpha},\overline{e^{\imath\alpha}})[X]$, the following is  equivalent: 
\begin{itemize}
    \item $\exp(\textbf{i} \alpha/3) \notin \mathcal{M}_{\infty}$
    \item $\text{degree}(\mu_{e^{\imath\alpha/3},K_0}) = 3$
    \item $X^3 - e^{\imath\alpha}$ is irreducible over $\mathbb{Q}(e^{\imath\alpha},\overline{e^{\imath\alpha}})$
\end{itemize}
The following section will demonstrate that the angle of $\frac{\pi}{3}=60°$  is not trisectable.

\begin{lemma}
    \label{lem:degree_K0_angel}
    \leanok
    \lean{degrre_root_three, eaual_adjion_root_three}
    The degree of $K_0 = \mathbb{Q}(e^{\imath\frac{\pi}{3}},\overline{e^{\imath\frac{\pi}{3}}})$ is equal  to $2$. 
\end{lemma}
\begin{proof}
    \leanok
    \uses{def:K_M_0}
    For all real numbers $\alpha$, we have that  $$ \exp(\imath \alpha) = \cos(\alpha) + \imath \sin(\alpha).$$
    For $\alpha = \pi / 3$ we get
    \begin{equation*}
        \cos(\alpha) = \frac{1}{2}\qquad \text{and}\qquad \sin(\alpha) = \frac{\sqrt{3}}{2}
    \end{equation*}
    Therefore $\mathbb{Q}(e^{\imath\frac{\pi}{3}},\overline{e^{\imath\frac{\pi}{3}}})$ is in $\mathbb{Q}(\imath\sqrt{3})$.
    And since $\imath\sqrt{3}$ is a zero of $X^2 + 3$, we know that the degree of $K_0$ less then $2$.
    To show that the degree is not $1$, we apply the fact that $\imath\sqrt{3} \notin \mathbb{Q}$.
    % Since we know that $\sqrt{r} \in \mathcal{M}_{\infty}$ for $r \in \mathcal{M}_n$
    % we see that $\exp(\textbf{i} \alpha) \in \mathcal{M}_{\infty}$ for $\mathcal{M} = \{0,1\}$. \newline
    % So we will work with $K_0 = \mathbb{Q}$.
\end{proof}

\begin{lemma}
    \label{lem:pi_third_not_in_M_inf}
    \leanok
    \lean{pi_third_not_in_M_inf}
    The angle $\pi / 3 = 60^{\circ}$ can't be trisected using a compass and straightedge.
\end{lemma}
\begin{proof}
    \leanok 
    \uses{lem:degree_K0_angel, cor:Classfication_z_in_M_inf_2m_rat}
    By utilising the aforementioned lemma \ref{lem:degree_K0_angel} to apply the corresponding corollary \ref{cor:Classfication_z_in_M_inf_2m_rat}, we can narrow our focus to the degree over $\Q$.
    Now we use the fact that if $x\in \mathcal{M}_{\infty}$, then $x.real, x.imag \in \mathcal{M}_{\infty}$ \ref{lem:M_inf_properties}. 
    Thus we focus on $\cos(\alpha/3)$, which the real part $e^{\imath \frac{\pi}{3}}$ and a zero of
    \begin{equation*}
        f := 8 X^3 - 6 X - 1 \in \mathbb{Q}[X]
    \end{equation*}
    Suppose $f$ is reducible over $\mathbb{Q}$, then $f$ has a rational zero $a$, since $f$ is of degree $3$. According to the rational root theorem, a root rational root of $f$ is of the form $\pm \frac{p}{q}$ with $p$ a factor of the constant term and $q$ a factor of the leading coefficient. So the only possible rational zeros of $f$ are
     \begin{equation*}
        \{ \pm 1, \pm \frac{1}{2}, \pm \frac{1}{4}, \pm \frac{1}{8} \}.
     \end{equation*}
     One can check that none of these numbers are a zero of $f$.
     So $f$ is irreducible over $\mathbb{Q}$ and $\cos(\alpha/3) \notin \mathcal{M}_{\infty}$.
     Therefore
        \begin{equation*}
            \exp(\textbf{i} \alpha/3) \notin \mathcal{M}_{\infty}
        \end{equation*}
    So the angle $\pi / 3 = 60^{\circ}$ can't be trisected using a compass and straightedge.
\end{proof}

\begin{theorem}
    \label{thm:Angle_not_Trisectable}
    \leanok
    \lean{Angle_not_Trisectable}
    A general angle can't be trisected using a compass and straightedge.
\end{theorem}
\begin{proof}
    \leanok
    \uses{lem:pi_third_not_in_M_inf, }
    Employ the previous lemma with the angle $\pi / 3$.
\end{proof}

%%FILE blueprint/src/chapter/chtopbestiary.tex
\chapter{Appendix: A collection of results which are needed in the proof.}
\label{ch_bestiary}
In this (temporary, unorganised) appendix I list a whole host of definitions and theorems which were known to humanity by the end of the univers \cite{JAN_SCHRÖER:2023}. These definitions and theorems will find their way into more relevant sections of the blueprint once I have time. Note also that many of the \emph{definitions} here are yet to be formalised in Lean, and this needs to be done before we can start talking about formalising theorems.
%We need the a debug placeholder for BibTex Stuff:


\section*{Stuff I'm working on}



\section{Kapitel 2 stuff for safekeeping}

%\begin{remark}
%    We know that $K_0$ is a field and $\mathcal{M} \subseteq K_0\subseteq \mathcal{M}_{\infty} \subseteq \mathbb{C}$.
%\end{remark}
\begin{theorem}
\label{degree_of_constructable_points}
    Let $\mathcal{M}\subseteq\mathbb{C}$ with $0,1 \in \mathcal{M}$ and $K_0 := \mathbb{Q}(\mathcal{M}\cup \overline{\mathcal{M}})$.
    Then for $z \in \mathcal{M}_{\infty}$ is equvalent:
    \begin{enumerate}
        \item $z \in \mathcal{M}_{\infty}$
        \item There is an $n \in \mathbb{N}$ and a chain \begin{equation*}
            K_0 = L_0 \subset L_1 \subset \dots \subset L_n \subset \mathbb{C}
        \end{equation*}
        of subfields of $\mathbb{C}$ such that $z \in L_n$ and $[L_i:L_{i-1}] =2$ for $1\le i\le n$ %$L_{i+1}$ is obtained from $L_i$ by adjoining a square root.
    \end{enumerate}
    In this case $[K_0(z):K_0] = 2^m$ for some $0 \le m \le n$.
\end{theorem}
\begin{remark}
    Theorem \ref{degree_of_constructable_points} tells us that it is sufficient to show that $[K_0(z):K_0] \ne 2^m$ for some $0 \le m $, witch we will use to show that $\sqrt[3]{2} \notin \mathcal{M}_{\infty}$ and $\exp(\textbf{i} \alpha/3) \notin \mathcal{M}_{\infty}$ for some $\alpha$.
\end{remark}
\begin{proof}[Proof of Theorem \ref{degree_of_constructable_points}]
    TODO %TODO
\end{proof}

\section{Safekeeping constructions}

%2.3.4
\begin{lemma}[Mitpiont of two complex numbers]
    \label{construction_midpoint}
    For $z_1, z_2 \in M_{\infty}$ is $\frac{z_1 + z_2}{2} \in M_{\infty}$.
\end{lemma}
%TODO: add proof, and Construction Idea

%2.3.5
\begin{lemma}[Halving of Angle]
    \label{lem:Construction.angle_half_M_inf}
    For $z = r \exp(\imath \alpha) \in M_{\infty}$ is $\exp(\imath \alpha / 2) \in M_{\infty}$.
\end{lemma}

%TODO: add proof, and Construction Idea

%2.3.6
%TODO: maybe write as \exp(\imath \alpha) + \exp(\imath \beta) = \exp(\imath \gamma) \in M_{\infty}
\begin{lemma}[Addition of Angles]
    \label{construction_add_angle}
    For two given angle $\alpha$ (given by two streches $[a,b]$ and $[c,d]$) and $\beta$ (given by two streches $[e,f]$ and $[g,h]$), with $ a,b,c,d,e,f,g,h \in M_{\infty}$, the addition of the angles is constructable.
\end{lemma}
%TODO: add proof, and Construction Idea, 

%2.3.7
\begin{lemma}[multiplication of positive real numbers]
    \label{construction_mult_pos_real}
    For $r_1, r_2 > 0 \in \mathbb{R}\cup M_{\infty}$ is $r_1 r_2 \in M_{\infty}$.
\end{lemma}
%TODO: add proof, and Construction Idea

%2.3.8
\begin{lemma}[Inverse of positive real numbers]
    \label{construction_inv_pos_real}
    For $r > 0\in \mathbb{R}\cup M_{\infty}$ is $r^{-1} \in M_{\infty}$.
\end{lemma}
%TODO: add proof, and Construction Idea

%2.3.9
\begin{lemma}[Root of positive real numbers]
    \label{construction_sqrt_pos_real}
    For $r > 0 \in \mathbb{R}\cup M_{\infty}$ is $\sqrt{r} \in M_{\infty}$.
\end{lemma}
%TODO: add proof, and Construction Idea

%2.3.10
\begin{lemma}[Polarkoordinates of complex numbers]
    \label{construction_polar}
    For $z = r \exp(\imath \alpha) \in M_{\infty}$ is $r, \allowbreak \exp(\imath \alpha) \in M_{\infty}$.
\end{lemma}
%TODO: add proof, and Construction Idea

%2.3.11
\begin{lemma}[Real and Imaginary part of complex numbers]
    \label{construction_re_im}
    For $0 \ne z = x + \imath y \in M_{\infty}$ is $x, \imath y \in M_{\infty}$.
\end{lemma}
%TODO: add proof, and Construction Idea

%2.3.12
\begin{lemma}[Construction of $\imath$]
    \label{construction_imath}
    $\imath \in M_{\infty}$.
\end{lemma}
%TODO: add proof, and Construction Idea

\section{Calssification}
Now we want to find a critirum for a element of $\mathbb{C}$ to be constructable, i.e. to be in $\mathcal{M}_{\infty}$.

\begin{theorem}
    \label{thm:Z_in_Minf_imp}
    For $z \in \mathbb{C}$. $z \in \mathcal{M}_{\infty}$ is equivalent to:\\
    There is an $n\le 0$ and chain 
    $$K_0 = L_1 \subset L_2 \subset \ldots \subset L_n \subset \mathbb{C}$$
    of subfields of $\mathbb{C}$ so that $z \in L_n$ and 
    $ [L_{i+1}:L_i] = 2 \quad \text{for} \quad i = 0, \ldots, n-1$.
\end{theorem}

\begin{lemma}
    \label{lem:Z_in_Minf_imp_eq} 
    For $z \in \mathbb{C}$. \\
    The exists of an $n\le 0$ and chain 
    $$K_0 = L_1 \subset L_2 \subset \ldots \subset L_n \subset \mathbb{C}$$
    of subfields of $\mathbb{C}$ so that $z \in L_n$ and 
    $ [L_{i+1}:L_i] = 2 \quad \text{for} \quad i = 0, \ldots, n-1$ is equivalent to:\\
    There is an intermidiate field $L$ of $\mathbb{C}/K_0$ whit $z \in L$, so that $L$ aries from $K_0$ by a sequence of adjunctions of square roots.
\end{lemma}

\begin{lemma}
    \label{lem:Mi_chain}
    For $M_n$ existiert a chain of intermidiate Fields $K_0 \le K_1 \le \ldots \le K_n$ so that $M_i\subset K_i$ and $K_i:= K_{i-1}(X_i)$ for a set of square roots $X_i$ of elements of $K_{i-1}$.
\end{lemma}

\begin{proof}
Induction over $n$

\end{proof}

\begin{lemma} 
    \label{lem:Z_in_Minf_imp_eq2}
    For $z \in \mathbb{C}$. $z \in \mathcal{M}_{\infty}$ is implies that there exists an $m \in \mathbb{N}$ such that $[z:K_0] = [K_0(z):K_0]= 2^m$.
\end{lemma}

\begin{proof} \uses{thm:Z_in_Minf_imp} %Todo add uses
    Let $z \in \mathbb{C}$ be an element of  $\mathcal{M}_{\infty}$.\\
    After theorem \ref{thm:Z_in_Minf_imp}, there exist an intermidiate field $L_n$ of $\Q$ $\C$ such that $z\in L_n$.\\
    Clame 1: $[L_n:K_0]= 2^n$\\

    Sinces ther is a chain \\
    Since $z\in L_n$ $K_0(z)$ is a gdsiufos $2^n=[L_n : K_0] = [L_n : K_0(z)] \cdot [K_0(Z):K_0]$.
    So $[K_0(Z):K_0] | 2^n$, since the only diversos of $2^n$ is $2^i$ for $\forall 0\le i \le n$ $[K_0(Z):K_0]$ is equal to $2^m$
\end{proof}

\begin{corollary}
    \label{cor:not_in_M_inf}
    For $z \in \mathbb{C}$, if ther is no $m$ such that $[z:K_0]=2^m$ then $z\notin \mathcal{M}_{\infty}$.
\end{corollary}

\begin{proof}
    \uses{lem:Z_in_Minf_imp_eq2}
    Follws diractly from \ref{lem:Z_in_Minf_imp_eq2}
\end{proof}

\section[Filed of Constructable Numbers]{Field $M_{\infty}$}
%2.3
\begin{theorem}
    % \label{thm:MField}
    % \leanok
    % \lean{MField}
    For $M\subseteq \mathbb{C}$ with $0,1 \in M$ is $M_{\infty}$ a subfield of $\mathbb{C}$.
\end{theorem}
\begin{proof}
    %TODO: uses M_inf M_M_inf add_M_Inf z_neg_M_inf mul_M_inf inv_M_inf
    %TODO: change proof from lean to latex 
    Chose $z_1, z_2 \in M_{\infty}$,after construction \ref{lem:construction_add}
    we have $z_1 + z_2 \in M_{\infty}$ and construction \ref{lem:construction_neg}
    gives us that for every $z \in M_{\infty}$ we have $-z \in M_{\infty}$. Therefore $M_{\infty}$ is closed under addition, and contains the additive invers.\\
    For multiplication let $z_1 = r_1 \exp(\imath \alpha_1)$ and $z_2 = r_2 \exp(\imath \alpha_2)$ be in $M_{\infty}$. Then we have 
    $$z_1 z_2 = r_1 r_2 \exp(\imath (\alpha_1 + \alpha_2)) \in M_{\infty}$$. By combining the lemmas \ref{construction_add_angle}, \ref{construction_mult_pos_real} and \ref{construction_polar}
    we obtain that $z_1 z_2 \in M_{\infty}$.\\
    For the inverse of $0 \ne z = r \exp(\imath \alpha) \in M_{\infty}$ we know that $$z^{-1} = r^{-1} \exp(\imath(2\pi - \alpha)).$$ We know that $r, \exp(\imath \alpha) \in M_{\infty}$, see construction \ref{construction_polar}.
    Now we can use construction \ref{construction_inv_pos_real} and \ref{lem:construction_conj} to show,
    that $r^{-1}, \exp(\imath(2\pi - \alpha)) = \overline{\exp(\imath\alpha)} \in M_{\infty}$. By applying that $M_{\infty}$ is closed under multiplication, shown above, we get that $z^{-1} \in M_{\infty}$.\\
    Therefore $M_{\infty}$ is a subfield of $\mathbb{C}$.
\end{proof}

\begin{corollary} 
    %TODO make to remark
    %\label{cor:MField_field}
    %\leanok
    %\lean{MField_field}
    For $M\subseteq \mathbb{C}$ with $0,1 \in M$ is $M_{\infty}$ a field.
\end{corollary}
\begin{proof}
    \uses{thm:MField}
    This is a direct consequence of Theorem \ref{thm:MField} and the fact that a subfield of a field is a field.
\end{proof}


%2.4
\begin{lemma}
    % \label{lem:M_inf_properties}
    % \leanok
    % \lean{MField_i, MField_re_im, MField_re_im', MField_polar}
    % \uses{thm:MField}
    For $M\subseteq \mathbb{C}$ with $0,1 \in M$ applies:
    \begin{itemize}
        \item[(i)] $\imath \in M_{\infty}$.
        \item[(ii)] For $z = x + \imath y \in \mathbb{C}$ is equivalent:
            \begin{itemize}
                \item $z \in M_{\infty}$.
                \item $x, y \in M_{\infty}$.
                \item $x, \imath y \in M_{\infty}$.
            \end{itemize}
        \item[(iii)] 
            For $0 \ne z = r \exp(\imath \alpha) \in \mathbb{C}$ is equivalent:
            \begin{itemize}
                \item $z \in M_{\infty}$.
                \item $r,\exp(\imath \alpha) \in M_{\infty}$.
            \end{itemize}
    \end{itemize}
\end{lemma}
\begin{proof}
    %TODO: \uses{} imath_M_inf z_iff_re_im_M_inf; (i_z_imp_z_in_M_inf/ MField_i_mul)
    This lemma is a dierct consequence of section \ref{basic_constructions}.
    \begin{itemize}
        \item[(i):] We kan apply construction \ref{construction_imath}
        \item[(ii):] We can apply construction \ref{construction_re_im} and \ref{construction_imath}.
        \item[(iii):]  We can apply construction \ref{construction_polar}.
    \end{itemize}
\end{proof}

\begin{definition}[quadritc closed field]
    % \label{def:quadritc_closed_field}
    % \leanok
    % \lean{QuadraticClosed}
    A field $K$ is called \emph{quadritc closed} if for all $z \in K$ we have $\sqrt{z} \in K$.
\end{definition}
\begin{remark}
    An equivalent definition is that $K$ is quadritc closed if $$K=\{a^2\mid a \in K\}.$$ 
\end{remark}

%2.6
\begin{lemma}
    %\label{lem:M_inf_quad_closed}
    %\leanok
    %\lean{MField_quad_closed}
    \uses{def:set_of_constructable_points, def:quadritc_closed_field}
    For $M\subseteq \mathbb{C}$ with $0,1 \in M$ is $M_{\infty}$ quadritc closed.
\end{lemma}
\begin{proof}
    \uses{cor:MField_field, cor:root_M_inf}
    Combian theorem \ref{cor:MField_field} to get that $M_{\infty}$ is a field and Lemma \ref{cor:root_M_inf} to get that $M_{\infty}$ is a quadritc closed field.
\end{proof}

\begin{definition}
    % \label{def:conjugate_set}
    % \leanok
    % \lean{conj_set}
    For a Set $U \subset \mathbb{C}$ we define the \emph{conjugate set} of $U$ as 
    \begin{equation*}
        Conj(U) = \{z\in \mathbb{C} \mid \exists w\in U: z = \overline{w}\}
    \end{equation*}
\end{definition}

\begin{definition}
    % \label{def:conj_closed}
    % \leanok
    % \lean{ConjClosed}
    We call a subset of $\C$ \emph{conjugate closed} if $M= Conj(M)$.
\end{definition}

\begin{lemma}
    % \label{lem:conj_MField}
    % \leanok
    % \lean{MField_conj}
    $M_{\infty}$ is conjugate closed.
\end{lemma}
\begin{proof}
    We can apply construction \ref{lem:construction_conj} and the fact that $\overline{\overline z} = z$ for all $z \in\mathbb{C}$.
\end{proof}




\section{Joshuas Kapitel}

\begin{definition}
    The degree of $x$ over $K$ is
    \begin{equation*}
        [x:K] :=\text{degree}(\mu_{x,K})
        \end{equation*}
        with $\mu_{x,K}$ the minimal polynomial of $x$ over $K$. \newline
    The degree of $L/K$ is the dimension of $L$ as a $K$-vector space and is denoted by
    \begin{equation*}
        [L:K].
    \end{equation*}
\end{definition}

\begin{theorem}
%\label{thm:degree_of_simple_field_extension}
    Let $L/K$ be a simple field extension with $L = K(x)$. Then
    \begin{equation*}
        [L:K] = [x:K].
    \end{equation*}
\end{theorem}
\begin{proof}
    In Mathlib: theorem IntermediateField.adjoin.finrank
\end{proof}

\begin{definition}
    %\label{def:K_M_0}
    Let $\mathcal(M)\subseteq\mathbb{C}$ with $0,1 \in \mathcal{M}$
    \begin{equation*}
        K_0 := \mathbb{Q}(\mathcal{M}\cup \overline{\mathcal{M}})
    \end{equation*}
    with $\overline{\mathcal{M}} := \{ \overline{z} = x - \textbf{i}y \mid z = x+\textbf{i}y  \in \mathcal{M} \}$.
\end{definition}

TODO: Add $K_i$ and $K_{\infty}$ and a Lemma that states that $K_{\infty} = M_{\infty}$. But this is not needed for the project.

%3.20
\begin{lemma}
    %\label{lem:conj_of_adjoin}
    Let $K$ be a subfield of $\mathbb{C}$ with $conj(K)=K$, and $ M \subset \mathbb{C}$ be a subset with $M = conj(M)$. Then holds
    \begin{equation*}
        conj(K(M)) = K(M).
    \end{equation*}
\end{lemma}
\begin{proof}
    We know that the complex conjugation is a field automorphism $\overline{\cdot }: \mathbb{C} \to \mathbb{C}$, whit $\overline{\overline{z}} = z$ for all $z \in \mathbb{C}$. Therefore we have
    $conj(K(M))$ is a subfield of $\mathbb{C}$.\\
    Since $K$ and $M$ are subsets of $K(m)$ we have $conj(K) = K$ and $conj(M) = M$ are subsets of $conj(K(M))$. Therefore $$K(M) \subset conj(K(M)).$$
    If we know apply $conj$ to both sides we get
    $$conj(K(M)) \subset conj(conj(K(M))) = K(M).$$
    Therefore $conj(K(M)) = K(M)$.
\end{proof}

%3.19
\begin{lemma}
    %\label{lem:dergee2_iff_adjoint_sqrt}
    Let $K\subset L \subset \mathbb{C}$ be fields. Then the following is equivalent:
    \begin{itemize}
        \item $[L:K] = 2$.
        \item There is a $w \in K$ with $\sqrt{w} \notin K$ and $L = K(\sqrt{w})$.
    \end{itemize}
\end{lemma}


%TODO: add lemma wich $K_0 \subset M_{\infty}$

%3.21
\begin{theorem}[constructable iff chain dergee2]
    %\label{thm:Z_in_Minf_1}
    For $z \in \mathbb{C}$ is equivalent:
    \begin{itemize}
        \item $z \in M_{\infty}$
        \item There is an intermidiate field $L$ of $\mathbb{C}/K_0$ whit $z \in L$, so that $L$ aries from $K_0$ by a sequence of adjunctions of square roots.
        \item There is an $n\le 0$ and chain 
            $$K_0 = L_1 \subset L_2 \subset \ldots \subset L_n \subset \mathbb{C}$$
            of subfields of $\mathbb{C}$ so that $z \in L_n$ and 
            $$ [L_{i+1}:L_i] = 2 \quad \text{for} \quad i = 0, \ldots, n-1.$$
    \end{itemize}
\end{theorem}

\begin{remark}
    %\label{rem:Z_in_Minf_1}
    In this cases the $[K_0(z):K_0] = 2^m$ for some $0 \le m \le n$.
\end{remark}
\begin{proof}[Proof of Theorem \ref{thm:Z_in_Minf_1}]
    \uses{ lem:dergee2_iff_adjoint_sqrt, Intersection_circle_circle, Intersection_line_circle, Intersection_line_line, lem:conj_of_adjoin, lem:M_inf_quad_closed}
    $(ii) \iff (iii)$ Follows from the definition and Lemma \ref{lem:dergee2_iff_adjoint_sqrt}.\\
    $(i) \implies (ii)$: We assume that $z$ is derived from $K_0$ by (ZL1), (ZL2) and (ZL3). %TODO: add ref
    Lemma \ref{Intersection_circle_circle}, \ref{Intersection_line_circle} and \ref{Intersection_line_line} in combination with Lemma \ref{lem:conj_of_adjoin} 
    imply that ther is an $w \in K_0$ so that $z \in K_0(\sqrt{w})$. Since $\overline{w} \in K_0$, by definition. We define
    $$K_1 := K_0(\sqrt{w},\sqrt{\overline{w}}).$$
    $K_1$ derives from $K_0$ sucsessively by adjunction of square roots. Note that $K_1 = conj(K_1)$, by Lemma \ref{lem:conj_of_adjoin} and the fact that $\sqrt{\overline{w}} = \overline{\sqrt{w}}$.
    Because of the definition of $$M_{\infty} = \bigcup_{n=0}^{\infty} M_n,$$
    he statement follows by induction.\\ %TODO: Elaborate on induction for lean
    $(ii) \implies (i)$: We know that $M_{\infty}$ is an quadritc complet field \ref{lem:M_inf_quad_closed}, i.e. $\sqrt{z} \in M_{\infty}$ for all $z \in M_{\infty}$.
    Futhermore we know that $K_0\subseteq M_{\infty} \subseteq \mathbb{C}$, %TODO: add ref
    therefore every field $L\subseteq \mathbb{C}$ which is derived from $K_0$ by a sequence of adjunctions of square roots is contaianted in $M_{\infty}$.
\end{proof}

%%FILE blueprint/src/content.tex
% In this file you should put the actual content of the blueprint.
% It will be used both by the web and the print version.
% It should *not* include the \begin{document}
%
% If you want to split the blueprint content into several files then
% the current file can be a simple sequence of \input. Otherwise It
% can start with a \section or \chapter for instance.
\input{chapter/ch01introduction}
\input{chapter/ch02fieldofconstnum}
\input{chapter/ch0xancientconstructionProplems}
%\input{chapter/chtopbestiary} -- The shloud be no longer needed

% Need to be moved for the print version
% \input{chapter/biblio}
% \nocite{*}

%%FILE blueprint/src/lstlean.tex
% Listing style definition for the Lean Theorem Prover.
% Defined by Jeremy Avigad, 2015, by modifying Assia Mahboubi's SSR style.
% Unicode replacements taken from Olivier Verdier's unixode.sty

\lstdefinelanguage{lean} {

% Anything between $ becomes LaTeX math mode
mathescape=false,
% Comments may or not include Latex commands
texcl=false,

% keywords, list taken from lean-syntax.el
morekeywords=[1]{
import, prelude, protected, private, noncomputable, definition, meta, renaming,
hiding, parameter, parameters, begin, constant, constants,
lemma, variable, variables, theory,
print, theorem, example,
open, as, export, override, axiom, axioms, inductive, with,
structure, record, universe, universes,
alias, help, precedence, reserve, declare_trace, add_key_equivalence,
match, infix, infixl, infixr, notation, postfix, prefix, instance,
eval, reduce, check, end, this,
using, using_well_founded, namespace, section,
attribute, local, set_option, extends, include, omit, class,
raw, replacing,
calc, have, show, suffices, by, in, at, let, forall, Pi, fun,
exists, if, dif, then, else, assume, obtain, from, register_simp_ext, unless, break, continue,
mutual, do, def, run_cmd, const,
partial, mut, where, macro, syntax, deriving,
return, try, catch, for, macro_rules, declare_syntax_cat, abbrev},

% Sorts
morekeywords=[2]{Sort, Type, Prop},

% tactics, list taken from lean-syntax.el
morekeywords=[3]{
assumption,
apply, intro, intros, allGoals,
generalize, clear, revert, done, exact,
refine, repeat, cases, rewrite, rw,
simp, simp_all, contradiction,
constructor, injection,
induction,
},

% modifiers, taken from lean-syntax.el
% note: 'otherkeywords' is needed because these use a different symbol.
% this command doesn't allow us to specify a number -- they are put with [1]
% otherkeywords={
% [persistent], [notation], [visible], [instance], [trans_instance],
% [class], [parsing-only], [coercion], [unfold_full], [constructor],
% [reducible], [irreducible], [semireducible], [quasireducible], [wf],
% [whnf], [multiple_instances], [none], [decl], [declaration],
% [relation], [symm], [subst], [refl], [trans], [simp], [congr], [unify],
% [backward], [forward], [no_pattern], [begin_end], [tactic], [abbreviation],
% [reducible], [unfold], [alias], [eqv], [intro], [intro!], [elim], [grinder],
% [localrefinfo], [recursor]
% },

% Various symbols
literate=
{α}{{\ensuremath{\mathrm{\alpha}}}}1
{β}{{\ensuremath{\mathrm{\beta}}}}1
{γ}{{\ensuremath{\mathrm{\gamma}}}}1
{δ}{{\ensuremath{\mathrm{\delta}}}}1
{ε}{{\ensuremath{\mathrm{\varepsilon}}}}1
{ζ}{{\ensuremath{\mathrm{\zeta}}}}1
{η}{{\ensuremath{\mathrm{\eta}}}}1
{θ}{{\ensuremath{\mathrm{\theta}}}}1
{ι}{{\ensuremath{\mathrm{\iota}}}}1
{κ}{{\ensuremath{\mathrm{\kappa}}}}1
{μ}{{\ensuremath{\mathrm{\mu}}}}1
{ν}{{\ensuremath{\mathrm{\nu}}}}1
{ξ}{{\ensuremath{\mathrm{\xi}}}}1
{π}{{\ensuremath{\mathrm{\mathnormal{\pi}}}}}1
{ρ}{{\ensuremath{\mathrm{\rho}}}}1
{σ}{{\ensuremath{\mathrm{\sigma}}}}1
{τ}{{\ensuremath{\mathrm{\tau}}}}1
{φ}{{\ensuremath{\mathrm{\varphi}}}}1
{χ}{{\ensuremath{\mathrm{\chi}}}}1
{ψ}{{\ensuremath{\mathrm{\psi}}}}1
{ω}{{\ensuremath{\mathrm{\omega}}}}1

{Γ}{{\ensuremath{\mathrm{\Gamma}}}}1
{Δ}{{\ensuremath{\mathrm{\Delta}}}}1
{Θ}{{\ensuremath{\mathrm{\Theta}}}}1
{Λ}{{\ensuremath{\mathrm{\Lambda}}}}1
{Σ}{{\ensuremath{\mathrm{\Sigma}}}}1
{Φ}{{\ensuremath{\mathrm{\Phi}}}}1
{Ξ}{{\ensuremath{\mathrm{\Xi}}}}1
{Ψ}{{\ensuremath{\mathrm{\Psi}}}}1
{Ω}{{\ensuremath{\mathrm{\Omega}}}}1

{ℵ}{{\ensuremath{\aleph}}}1

{≤}{{\ensuremath{\leq}}}1
{≥}{{\ensuremath{\geq}}}1
{≠}{{\ensuremath{\neq}}}1
{≈}{{\ensuremath{\approx}}}1
{≡}{{\ensuremath{\equiv}}}1
{≃}{{\ensuremath{\simeq}}}1

{≤}{{\ensuremath{\leq}}}1
{≥}{{\ensuremath{\geq}}}1

{∂}{{\ensuremath{\partial}}}1
{∆}{{\ensuremath{\triangle}}}1 % or \laplace?

{∫}{{\ensuremath{\int}}}1
{∑}{{\ensuremath{\mathrm{\Sigma}}}}1
{Π}{{\ensuremath{\mathrm{\Pi}}}}1

{⊥}{{\ensuremath{\perp}}}1
{∞}{{\ensuremath{\infty}}}1
{∂}{{\ensuremath{\partial}}}1

{∓}{{\ensuremath{\mp}}}1
{±}{{\ensuremath{\pm}}}1
{×}{{\ensuremath{\times}}}1

{⊕}{{\ensuremath{\oplus}}}1
{⊗}{{\ensuremath{\otimes}}}1
{⊞}{{\ensuremath{\boxplus}}}1

{∇}{{\ensuremath{\nabla}}}1
{√}{{\ensuremath{\sqrt}}}1

{⬝}{{\ensuremath{\cdot}}}1
{•}{{\ensuremath{\cdot}}}1
{∘}{{\ensuremath{\circ}}}1

%{⁻}{{\ensuremath{^{\textup{\kern1pt\rule{2pt}{0.3pt}\kern-1pt}}}}}1
{⁻}{{\ensuremath{^{-}}}}1
{▸}{{\ensuremath{\blacktriangleright}}}1

{∧}{{\ensuremath{\wedge}}}1
{∨}{{\ensuremath{\vee}}}1
{¬}{{\ensuremath{\neg}}}1
{⊢}{{\ensuremath{\vdash}}}1

%{⟨}{{\ensuremath{\left\langle}}}1
%{⟩}{{\ensuremath{\right\rangle}}}1
{⟨}{{\ensuremath{\langle}}}1
{⟩}{{\ensuremath{\rangle}}}1

{↦}{{\ensuremath{\mapsto}}}1
{←}{{\ensuremath{\leftarrow}}}1
{<-}{{\ensuremath{\leftarrow}}}1
{→}{{\ensuremath{\rightarrow}}}1
{↔}{{\ensuremath{\leftrightarrow}}}1
{⇒}{{\ensuremath{\Rightarrow}}}1
{⟹}{{\ensuremath{\Longrightarrow}}}1
{⇐}{{\ensuremath{\Leftarrow}}}1
{⟸}{{\ensuremath{\Longleftarrow}}}1

{ᶠ}{{\ensuremath{^f}}}1
{∩}{{\ensuremath{\cap}}}1
{∪}{{\ensuremath{\cup}}}1
{⋃}{{\ensuremath{\bigcup}}}1
{⊂}{{\ensuremath{\subseteq}}}1
{⊆}{{\ensuremath{\subseteq}}}1
{⊄}{{\ensuremath{\nsubseteq}}}1
{⊈}{{\ensuremath{\nsubseteq}}}1
{⊃}{{\ensuremath{\supseteq}}}1
{⊇}{{\ensuremath{\supseteq}}}1
{⊅}{{\ensuremath{\nsupseteq}}}1
{⊉}{{\ensuremath{\nsupseteq}}}1
{∈}{{\ensuremath{\in}}}1
{∉}{{\ensuremath{\notin}}}1
{∋}{{\ensuremath{\ni}}}1
{∌}{{\ensuremath{\notni}}}1
{∅}{{\ensuremath{\emptyset}}}1

{∖}{{\ensuremath{\setminus}}}1
{†}{{\ensuremath{\dag}}}1

{ℕ}{{\ensuremath{\mathbb{N}}}}1
{ℤ}{{\ensuremath{\mathbb{Z}}}}1
{ℝ}{{\ensuremath{\mathbb{R}}}}1
{ℚ}{{\ensuremath{\mathbb{Q}}}}1
{ℂ}{{\ensuremath{\mathbb{C}}}}1
{⌞}{{\ensuremath{\llcorner}}}1
{⌟}{{\ensuremath{\lrcorner}}}1
{⦃}{{\ensuremath{\{\!|}}}1
{⦄}{{\ensuremath{|\!\}}}}1

{‖}{{\ensuremath{\|}}}1
{₁}{{\ensuremath{_1}}}1
{₂}{{\ensuremath{_2}}}1
{₃}{{\ensuremath{_3}}}1
{₄}{{\ensuremath{_4}}}1
{₅}{{\ensuremath{_5}}}1
{₆}{{\ensuremath{_6}}}1
{₇}{{\ensuremath{_7}}}1
{₈}{{\ensuremath{_8}}}1
{₉}{{\ensuremath{_9}}}1
{₀}{{\ensuremath{_0}}}1
{ᵢ}{{\ensuremath{_i}}}1
{ⱼ}{{\ensuremath{_j}}}1
{ₐ}{{\ensuremath{_a}}}1

{¹}{{\ensuremath{^1}}}1

{ₙ}{{\ensuremath{_n}}}1
{ₘ}{{\ensuremath{_m}}}1
{ₚ}{{\ensuremath{_p}}}1
{↑}{{\ensuremath{\uparrow}}}1
{↓}{{\ensuremath{\downarrow}}}1

{...}{{\ensuremath{\ldots}}}1
{·}{{\ensuremath{\cdot}}}1

{▸}{{\ensuremath{\triangleright}}}1

{Σ}{{\color{symbolcolor}\ensuremath{\Sigma}}}1
{Π}{{\color{symbolcolor}\ensuremath{\Pi}}}1
{∀}{{\color{symbolcolor}\ensuremath{\forall}}}1
{∃}{{\color{symbolcolor}\ensuremath{\exists}}}1
{λ}{{\color{symbolcolor}\ensuremath{\mathrm{\lambda}}}}1
{\$}{{\color{symbolcolor}\$}}1

{:=}{{\color{symbolcolor}:=}}1
{=}{{\color{symbolcolor}=}}1
{<|>}{{\color{symbolcolor}<|>}}1
{<\$>}{{\color{symbolcolor}<\$>}}1
{+}{{\color{symbolcolor}+}}1
{*}{{\color{symbolcolor}*}}1,

% Comments
%comment=[s][\itshape \color{commentcolor}]{/-}{-/},
morecomment=[s][\color{commentcolor}]{/-}{-/},
morecomment=[l][\itshape \color{commentcolor}]{--},

% Spaces are not displayed as a special character
showstringspaces=false,

% keep spaces
keepspaces=true,

% String delimiters
morestring=[b]",
morestring=[d]’,

% Size of tabulations
tabsize=3,

% Enables ASCII chars 128 to 255
extendedchars=false,

% Case sensitivity
sensitive=true,

% Automatic breaking of long lines
breaklines=true,
breakatwhitespace=true,

% Default style fors listingsred
basicstyle=\ttfamily\small,

% Position of captions is bottom
captionpos=b,

% Full flexible columns
columns=[l]fullflexible,


% Style for (listings') identifiers
identifierstyle={\ttfamily\color{black}},
% Note : highlighting of Coq identifiers is done through a new
% delimiter definition through an lstset at the beginning of the
% document. Don't know how to do better.

% Style for declaration keywords
keywordstyle=[1]{\ttfamily\color{keywordcolor}},

% Style for sorts
keywordstyle=[2]{\ttfamily\color{sortcolor}},

% Style for tactics keywords
keywordstyle=[3]{\ttfamily\color{tacticcolor}},

% Style for attributes
keywordstyle=[4]{\ttfamily\color{attributecolor}},

% Style for strings
stringstyle=\ttfamily,

% Style for comments
commentstyle={\ttfamily\footnotesize },

}
